% AI for All, Non-Tech: Part A (non-tech core, ai_core_nontech_*) + Part B (everyday life, fluency, ChatGPT, applications).
% Pattern for every AI-for-X seminar: core first, then the audience part, then wrap-up. See AI_FOR_X_PLAN.md.
% To shorten, comment frames in place (do not delete them).

\section[Core]{AI Basics}
% Part A (non-tech core) of the AI-for-X seminars: what AI and generative AI are. 4 frames.
% Adapted from ai_educators_content.tex (Big Picture, Language Model, Large). Audience-neutral on purpose.
% Part A is shared: edit here, then recompile every deck that includes it. Trim by commenting frames, not deleting.
% Size target: same as ai_core_tech_whatisai (4 frames).

%%%%%%%%%%%%%%%%%%%%%%%%%%%%%%%%%%%%%%%%%%%%%%%%%%%%%%%%%%%
\begin{frame}[fragile]\frametitle{The Big Picture: AI, ML, and Deep Learning}
\begin{center}
\includegraphics[width=0.55\linewidth,keepaspectratio]{ai1}
\end{center}
\begin{itemize}
\item \textbf{AI}: any machine doing something that looks intelligent
\item \textbf{Machine Learning (ML)}: a machine that learns the pattern from data, instead of being told the rule directly
\item \textbf{Deep Learning (DL)}: ML using many-layered neural networks, powerful enough to learn something as complex as language
\end{itemize}
{\tiny (Ref: https://blogs.nvidia.com/blog/2016/07/29/whats-difference-artificial-intelligence-machine-learning-deep-learning-ai/)}
\end{frame}

%%%%%%%%%%%%%%%%%%%%%%%%%%%%%%%%%%%%%%%%%%%%%%%%%%%%%%%%%%%
\begin{frame}[fragile]\frametitle{What Is a Language Model?}
\begin{itemize}
\item While typing an SMS, have you noticed it suggests the next word?
\item Those suggestions are not random: the phone has learned what word usually comes next from millions of past sentences.
\item Do this once, word by word, and you get a suggestion. Do it again and again, and you get a full sentence.
\item That is a \textbf{Language Model}, predicting the next word. Chained together, it writes whole paragraphs: a \textbf{Generative Model}.
\end{itemize}
\begin{center}
\includegraphics[width=0.4\linewidth,keepaspectratio]{chatgpt34}
\end{center}
\end{frame}

%%%%%%%%%%%%%%%%%%%%%%%%%%%%%%%%%%%%%%%%%%%%%%%%%%%%%%%%%%%
\begin{frame}[fragile]\frametitle{Why Are They Called ``Large''?}
\begin{itemize}
\item ChatGPT works on exactly this idea, just at an enormous scale: trained on a large fraction of everything ever written on the internet.
\item That is why the phrase is \textbf{Large Language Model (LLM)}: large training data, large number of internal parameters.
\item More data and more parameters mean better predictions, but the underlying idea never changes: predict the next word, one at a time.
\end{itemize}
\begin{center}
\includegraphics[width=0.5\linewidth,keepaspectratio]{chatgpt31}
\end{center}
\end{frame}

%%%%%%%%%%%%%%%%%%%%%%%%%%%%%%%%%%%%%%%%%%%%%%%%%%%%%%%%%%%
\begin{frame}[fragile]\frametitle{Generative AI: What It Can Make}
\begin{itemize}
\item \textbf{Generative AI} makes new content from a request: text, images, code, audio, even short video.
\item The chat assistants you may have met (ChatGPT, Claude, Gemini and others) are built on language models like the one on the last slide.
\item Same idea, different material: the model has seen very many examples of the thing, and produces a new one that looks like them.
\item It is a very fast, very well-read assistant. It is not a person, and it does not understand your situation unless you tell it.
\end{itemize}
\end{frame}

% Part A (non-tech core) of the AI-for-X seminars: limits and responsible use. 4 frames.
% Adapted from ai_educators_wrapup.tex (failure modes, data) and ai_intro_nontech.tex. Audience-neutral on purpose.
% Part A is shared: edit here, then recompile every deck that includes it. Trim by commenting frames, not deleting.
% Size target: same as ai_core_tech_limits (4 frames).

%%%%%%%%%%%%%%%%%%%%%%%%%%%%%%%%%%%%%%%%%%%%%%%%%%%%%%%%%%%
\begin{frame}[fragile]\frametitle{It Predicts, It Does Not Know}
\begin{itemize}
\item A language model writes the most likely next words. It does not look facts up, and it does not know when it is wrong.
\item So it can produce a fluent, confident answer that is false: a \textbf{hallucination}. Invented names, dates, quotations and references are the classic cases.
\item The more your question depends on one exact fact, the more you must check the answer.
\end{itemize}
\begin{alertblock}{Rule of Thumb}
Fluent does not mean true. Check every fact, name and reference before you rely on it.
\end{alertblock}
\end{frame}

%%%%%%%%%%%%%%%%%%%%%%%%%%%%%%%%%%%%%%%%%%%%%%%%%%%%%%%%%%%
\begin{frame}[fragile]\frametitle{Other Limits Worth Knowing}
\begin{itemize}
\item \textbf{Out of date:} the model learned from text up to some date. Recent events need a web search or a source you supply.
\item \textbf{Bias:} models trained on internet-scale text inherit internet-scale biases. Read the output with that in mind.
\item \textbf{Inconsistency:} ask the same question twice and the answers can differ. Do not treat one answer as a ruling.
\item \textbf{Over-reliance:} a first draft is easy to accept as the final answer. It is a starting point, not a verdict.
\end{itemize}
\end{frame}

%%%%%%%%%%%%%%%%%%%%%%%%%%%%%%%%%%%%%%%%%%%%%%%%%%%%%%%%%%%
\begin{frame}[fragile]\frametitle{What Not to Paste}
\begin{itemize}
\item Public AI tools may keep what you type, and may use it to improve their models, depending on the tool and its settings.
\item \textbf{Never paste} personal data, passwords, contracts, unpublished work or anything confidential into a public AI tool.
\item If you need help with sensitive text, remove the names and details first, or use a tool your organisation has approved.
\end{itemize}
\begin{alertblock}{Rule of Thumb}
If you would not leave it on a sticky note on a public desk, do not paste it into a public AI chat.
\end{alertblock}
\end{frame}

%%%%%%%%%%%%%%%%%%%%%%%%%%%%%%%%%%%%%%%%%%%%%%%%%%%%%%%%%%%
\begin{frame}[fragile]\frametitle{Treat the Output as a First Draft}
\begin{itemize}
\item Read it, correct it, and make it yours before it goes anywhere.
\item You are responsible for what you publish, send or submit, whoever or whatever produced the first version.
\item Ask the tool to show its reasoning or sources, then check them yourself.
\item When in doubt, ask a person who knows the subject.
\end{itemize}
\end{frame}


\section[Prompts]{Prompts with CRAFT}
% Part A (non-tech core) of the AI-for-X seminars: prompts and the CRAFT framework. 8 frames.
% Adapted from llm_promptengg_educators.tex (What Is a Prompt, Same Task, CRAFT x5, Takeaways) with neutral
% examples. The Educators deck keeps its own education examples in its Part B.
% Part A is shared: edit here, then recompile every deck that includes it. Trim by commenting frames, not deleting.
% Size target: same as ai_core_tech_prompts (8 frames).

%%%%%%%%%%%%%%%%%%%%%%%%%%%%%%%%%%%%%%%%%%%%%%%%%%%%%%%%%%%
\begin{frame}[fragile]\frametitle{What Is a Prompt?}
\begin{itemize}
\item A \textbf{prompt} is the text you type to the AI: a question, an instruction, or both.
\item The model reads your prompt and continues it, one word at a time, just like the SMS suggestion.
\item So your words steer the answer: more specific text gives a more specific result.
\item A prompt can hold a role for the AI, background, an example, and the text to work on.
\item \textbf{Prompt engineering} is writing prompts that give usable answers.
\end{itemize}
\end{frame}

%%%%%%%%%%%%%%%%%%%%%%%%%%%%%%%%%%%%%%%%%%%%%%%%%%%
\begin{frame}[fragile]\frametitle{Same Task, Different Prompts}
\begin{itemize}
\item Same AI, same topic: a vague prompt gives a generic answer, a specific one gives a usable answer.
\item A chat remembers what came before. Reply to refine the answer instead of starting again.
\end{itemize}
\begin{lstlisting}[language={}]
Vague:      Write a notice about water tank cleaning.

Specific:   Write a notice for a housing society about
            water tank cleaning on Sunday, 6 to 10 am,
            in 80 words, polite tone.

Follow-up:  Make it shorter, and add a contact number
            placeholder.
\end{lstlisting}
\end{frame}

%%%%%%%%%%%%%%%%%%%%%%%%%%%%%%%%%%%%%%%%%%%%%%%%%%%%%%%%%%%
\begin{frame}[fragile]\frametitle{A Framework for Every Prompt: CRAFT}

\begin{center}
\Large \textbf{C R A F T}
\end{center}

\medskip
\begin{itemize}
    \item \textbf{C}ontext, who you are, what the situation is
    \item \textbf{R}ole, who the AI should act as
    \item \textbf{A}ction, exactly what you want it to do
    \item \textbf{F}ewshots, one or two examples of what ``good'' looks like
    \item \textbf{T}weaks, tone, length, format, and what to avoid
\end{itemize}

\medskip
\begin{center}
\textbf{Run through the five letters in your head before you type. That is the whole framework.}
\end{center}

\end{frame}

%%%%%%%%%%%%%%%%%%%%%%%%%%%%%%%%%%%%%%%%%%%%%%%%%%%%%%%%%%%
\begin{frame}[fragile]\frametitle{CRAFT: Context and Role}

\begin{itemize}
    \item \textbf{Context}, without it, the AI is guessing: what the subject is, who the reader is, what has already happened.
    \item \textbf{Role}, telling the AI who to act as changes its whole register: a ``friendly neighbour'' writes differently from a ``formal clerk.''
\end{itemize}

\begin{lstlisting}[language=HTML]
Weak:  Write a notice about water tank cleaning.

CRAFT: You are the secretary of a housing society (Role).
       The tank is cleaned twice a year, and residents
       have complained about short notice before (Context).
\end{lstlisting}

\end{frame}

%%%%%%%%%%%%%%%%%%%%%%%%%%%%%%%%%%%%%%%%%%%%%%%%%%%%%%%%%%%
\begin{frame}[fragile]\frametitle{CRAFT: Action and Fewshots}

\begin{itemize}
    \item \textbf{Action}, one clear instruction. ``Help me with X'' is not an action, ``Write a 3-line notice about X'' is.
    \item \textbf{Fewshots}, show, do not just tell. One example of the style or format you want is worth three sentences of description.
\end{itemize}

\begin{lstlisting}[language=HTML]
Action:   Write a 3-line notice about the tank
          cleaning on Sunday.

Fewshot:  Match this style:
          "Water supply will be off on Monday,
           9 to 12. Please store water in advance."
\end{lstlisting}

\end{frame}

%%%%%%%%%%%%%%%%%%%%%%%%%%%%%%%%%%%%%%%%%%%%%%%%%%%%%%%%%%%
\begin{frame}[fragile]\frametitle{CRAFT: Tweaks}

\begin{itemize}
    \item \textbf{Tweaks} are the small dials that turn a generic answer into a usable one:
    \begin{itemize}
        \item \textbf{Tone}, formal / friendly / encouraging
        \item \textbf{Length}, word count or ``under 5 sentences''
        \item \textbf{Format}, bullet list / table / email / paragraph
        \item \textbf{Constraints}, what to explicitly avoid
    \end{itemize}
    \item Add tweaks last, once Context, Role, Action, and Fewshots already give the AI enough to work with.
\end{itemize}

\begin{lstlisting}[language=HTML]
Tweaks: Keep it under 60 words, friendly tone,
        no exclamation marks.
\end{lstlisting}

\end{frame}

%%%%%%%%%%%%%%%%%%%%%%%%%%%%%%%%%%%%%%%%%%%%%%%%%%%%%%%%%%%
\begin{frame}[fragile]\frametitle{CRAFT, All Together}

\begin{lstlisting}[language=HTML]
You are the secretary of a housing society (R).
The tank is cleaned twice a year, and residents
have complained about short notice before (C).

Write a 3-line notice about the tank cleaning
on Sunday, 6 to 10 am (A).

Match this style: "Water supply will be off on
Monday, 9 to 12. Please store water in advance." (F)

Keep it under 60 words, friendly tone,
no exclamation marks (T).
\end{lstlisting}

\end{frame}

%%%%%%%%%%%%%%%%%%%%%%%%%%%%%%%%%%%%%%%%%%%%%%%%%%%%%%%%%%%
\begin{frame}[fragile]\frametitle{CRAFT: The Takeaways}

\begin{itemize}
    \item Not every prompt needs all five letters spelled out explicitly, but running through them in your head catches what you would otherwise forget.
    \item Save a prompt that works. A small personal library of tested prompts saves more time than any single clever prompt.
    \item \textbf{Always remember:} treat the output as a first draft, review it, and never paste confidential data into a public AI tool.
\end{itemize}

\end{frame}


\section[Everyday]{AI in Everyday Life}
\begin{frame}[fragile]\frametitle{}
\begin{center}
{\Large AI in Everyday Life}
\end{center}
\end{frame}
%%%%%%%%%%%%%%%%%%%%%%%%%%%%%%%%%%%%%%%%%%%%%%%%%%%%%%%%%%%%%%%%%%%%%%%%%%%%%%%%%%
% [ai-for-x: commented 2026-11-07 | covered by the Part B divider "AI in Everyday Life" in seminar_ai_for_all_nontech_content.tex and ai_core_nontech_whatisai]
% \begin{frame}[fragile]\frametitle{}
% \begin{center}
% {\Large Understanding Artificial Intelligence}
%
% \vspace{1cm}
% {\large A Simple Guide for Everyone}
% \end{center}
% \end{frame}

%%%%%%%%%%%%%%%%%%%%%%%%%%%%%%%%%%%%%%%%%%%%%%%%%%%%%%%%%%%%%%%%%%%%%%%%%%%%%%%%%%
% [ai-for-x: commented 2026-11-07 | covered by the Part B divider "AI in Everyday Life" in seminar_ai_for_all_nontech_content.tex]
% \begin{frame}[fragile]\frametitle{}
% \begin{center}
% {\Large AI is Already Part of Your Daily Life!}
% \end{center}
% \end{frame}

%%%%%%%%%%%%%%%%%%%%%%%%%%%%%%%%%%%%%%%%%%%%%%%%%%%%%%%%%%%%%%%%%%%%%%%%%%%%%%%%%%
% [ai-for-x: commented 2026-11-07 | covered by ai_core_nontech_whatisai.tex]
% \begin{frame}[fragile]\frametitle{What is Artificial Intelligence?}
% Think of AI as:
% \begin{itemize}
% \item Teaching machines to be smart
% \item Like training a very clever assistant
% \item That never gets tired
% \item And can remember everything
% \item But needs YOUR guidance to work properly
% \end{itemize}
%
% \vspace{0.5cm}
% {\bf In simple words:} Making computers do tasks that normally need human intelligence
% \end{frame}

%%%%%%%%%%%%%%%%%%%%%%%%%%%%%%%%%%%%%%%%%%%%%%%%%%%%%%%%%%%%%%%%%%%%%%%%%%%%%%%%%%
\begin{frame}[fragile]\frametitle{AI in Your Everyday Life}

You're already using AI without realizing it!

\begin{itemize}
\item {\bf WhatsApp:} Suggests quick replies
\item {\bf Google Photos:} Groups photos by faces automatically
\item {\bf Netflix/Prime:} Recommends shows you might like
\item {\bf Gmail:} Filters spam emails
\item {\bf Facebook:} Shows posts you're likely to enjoy
\item {\bf Google Maps:} Predicts traffic and suggests routes
\item {\bf Amazon:} ``People who bought this also bought...''
\item {\bf YouTube:} Recommends next video to watch
\end{itemize}
\end{frame}

%%%%%%%%%%%%%%%%%%%%%%%%%%%%%%%%%%%%%%%%%%%%%%%%%%%%%%%%%%%%%%%%%%%%%%%%%%%%%%%%%%
\begin{frame}[fragile]\frametitle{Banking and Shopping}
\begin{itemize}
\item {\bf ATM Check Deposit:} Reads your check automatically
\item {\bf Credit Card:} Detects unusual purchases (fraud protection)
\item {\bf Online Shopping:} Personalized product recommendations
\item {\bf Virtual Assistants:} Alexa orders groceries by voice
\item {\bf Loan Applications:} Faster approval decisions
\end{itemize}

\vspace{0.5cm}
All powered by AI working behind the scenes!
\end{frame}

%%%%%%%%%%%%%%%%%%%%%%%%%%%%%%%%%%%%%%%%%%%%%%%%%%%%%%%%%%%%%%%%%%%%%%%%%%%%%%%%%%
\begin{frame}[fragile]\frametitle{}
\begin{center}
{\Large How Does This ``Intelligence'' Actually Work?}
\end{center}
\end{frame}

%%%%%%%%%%%%%%%%%%%%%%%%%%%%%%%%%%%%%%%%%%%%%%%%%%%%%%%%%%%%%%%%%%%%%%%%%%%%%%%%%%
\begin{frame}[fragile]\frametitle{Learning from Experience - Like Humans}

{\bf How you learned to cook:}
\begin{itemize}
\item First time: Maybe too salty or burnt
\item Next time: Adjusted based on what went wrong
\item After many attempts: You know exactly what to do
\item Now: You can even experiment with new recipes
\end{itemize}

{\bf AI learns the SAME way:}
\begin{itemize}
\item Give it many examples
\item It learns patterns
\item Gets better with practice
\item Can then handle new situations
\end{itemize}
\end{frame}

%%%%%%%%%%%%%%%%%%%%%%%%%%%%%%%%%%%%%%%%%%%%%%%%%%%%%%%%%%%%%%%%%%%%%%%%%%%%%%%%%%
\begin{frame}[fragile]\frametitle{Real Example: Email Spam Filter}

{\bf Old Way (Before AI):}
\begin{itemize}
\item Engineers write rules: ``If email says FREE MONEY, it's spam''
\item Spammers change words to ``FR33 M0NEY''
\item Rules break, spam gets through
\end{itemize}

{\bf New Way (With AI):}
\begin{itemize}
\item Show AI 10,000 spam emails
\item Show AI 10,000 real emails
\item AI learns subtle patterns (not just keywords!)
\item Adapts automatically when spammers change tactics
\end{itemize}

This is why your Gmail spam filter is so good!
\end{frame}

%%%%%%%%%%%%%%%%%%%%%%%%%%%%%%%%%%%%%%%%%%%%%%%%%%%%%%%%%%%%%%%%%%%%%%%%%%%%%%%%%%
\begin{frame}[fragile]\frametitle{Another Example: Voice Assistants}

When you say: {\em ``Alexa, what's the weather today?''}

\begin{itemize}
\item AI recognizes your voice (among all household members)
\item Converts speech to text
\item Understands you mean TODAY's weather, not tomorrow
\item Knows YOUR location automatically
\item Speaks response in natural language
\end{itemize}

All this happens in 2 seconds! Amazing, right?
\end{frame}

%%%%%%%%%%%%%%%%%%%%%%%%%%%%%%%%%%%%%%%%%%%%%%%%%%%%%%%%%%%%%%%%%%%%%%%%%%%%%%%%%%
% [ai-for-x: commented 2026-11-07 | covered by ai_core_nontech_whatisai.tex]
% \begin{frame}[fragile]\frametitle{}
% \begin{center}
% {\Large The AI Family - Different Types}
% \end{center}
% \end{frame}

%%%%%%%%%%%%%%%%%%%%%%%%%%%%%%%%%%%%%%%%%%%%%%%%%%%%%%%%%%%%%%%%%%%%%%%%%%%%%%%%%%
% [ai-for-x: commented 2026-11-07 | covered by ai_core_nontech_whatisai.tex]
% \begin{frame}[fragile]\frametitle{Understanding the AI Family}
%
% Think of it like vehicles:
%
% \begin{itemize}
% \item {\bf Artificial Intelligence:} All vehicles (the big category)
% \item {\bf Machine Learning:} Automatic transmission cars
% \begin{itemize}
% \item Learns from data automatically
% \end{itemize}
% \item {\bf Deep Learning:} Self-driving cars
% \begin{itemize}
% \item Very sophisticated, handles complex tasks
% \end{itemize}
% \item {\bf Generative AI:} Cars that can build new roads
% \begin{itemize}
% \item Creates NEW content: text, images, music
% \end{itemize}
% \end{itemize}
%
% Each is more advanced than the previous!
% \end{frame}

%%%%%%%%%%%%%%%%%%%%%%%%%%%%%%%%%%%%%%%%%%%%%%%%%%%%%%%%%%%%%%%%%%%%%%%%%%%%%%%%%%
\begin{frame}[fragile]\frametitle{ChatGPT - You May Have Heard About It}

{\bf What is ChatGPT?}
\begin{itemize}
\item An AI you can chat with, like texting a friend
\item Can answer questions, write emails, explain topics
\item Remembers your conversation
\item Learned from millions of books and websites
\end{itemize}

{\bf What people use it for:}
\begin{itemize}
\item Writing job applications
\item Planning meals with ingredients at home
\item Getting gift ideas
\item Explaining complex topics simply
\item Helping with homework (students!)
\end{itemize}
\end{frame}

%%%%%%%%%%%%%%%%%%%%%%%%%%%%%%%%%%%%%%%%%%%%%%%%%%%%%%%%%%%%%%%%%%%%%%%%%%%%%%%%%%
\begin{frame}[fragile]\frametitle{}
\begin{center}
{\Large AI in Different Aspects of Life}
\end{center}
\end{frame}

%%%%%%%%%%%%%%%%%%%%%%%%%%%%%%%%%%%%%%%%%%%%%%%%%%%%%%%%%%%%%%%%%%%%%%%%%%%%%%%%%%
\begin{frame}[fragile]\frametitle{AI in Healthcare}

{\bf For Patients:}
\begin{itemize}
\item Apps that monitor your health (heart rate, sleep)
\item Reminders to take medicines
\item Symptom checkers (though always see a real doctor!)
\end{itemize}

{\bf For Doctors:}
\begin{itemize}
\item AI reads X-rays and scans faster
\item Sometimes spots diseases doctors might miss
\item Suggests treatment options based on millions of cases
\end{itemize}

{\bf Important:} AI assists doctors, doesn't replace them!
\end{frame}

%%%%%%%%%%%%%%%%%%%%%%%%%%%%%%%%%%%%%%%%%%%%%%%%%%%%%%%%%%%%%%%%%%%%%%%%%%%%%%%%%%
\begin{frame}[fragile]\frametitle{AI in Home and Kitchen}

\begin{itemize}
\item {\bf Smart Thermostats:} Learn when you're home, adjust temperature
\item {\bf Smart Lights:} Turn on/off based on your routine
\item {\bf Recipe Apps:} Suggest dishes based on ingredients you have
\item {\bf Smart Ovens:} Recognize food, set cooking time automatically
\item {\bf Robotic Vacuums:} Clean house while you're away
\item {\bf Smart Fridges:} Track expiry dates, suggest recipes
\end{itemize}

Making life easier and more convenient!
\end{frame}

%%%%%%%%%%%%%%%%%%%%%%%%%%%%%%%%%%%%%%%%%%%%%%%%%%%%%%%%%%%%%%%%%%%%%%%%%%%%%%%%%%
\begin{frame}[fragile]\frametitle{AI in Transportation}

\begin{itemize}
\item {\bf Uber/Ola:} Matches you with nearest driver
\item {\bf Google Maps:} Predicts traffic, suggests faster routes
\item {\bf Flight Booking:} Predicts when prices will go up/down
\item {\bf Metro/Bus Apps:} Real-time arrival predictions
\item {\bf Parking Apps:} Find empty parking spots
\end{itemize}

{\bf Coming Soon:}
\begin{itemize}
\item Self-driving cars (already testing in some countries!)
\end{itemize}
\end{frame}

%%%%%%%%%%%%%%%%%%%%%%%%%%%%%%%%%%%%%%%%%%%%%%%%%%%%%%%%%%%%%%%%%%%%%%%%%%%%%%%%%%
\begin{frame}[fragile]\frametitle{AI in Entertainment}

\begin{itemize}
\item {\bf Music Apps:} Create playlists matching your mood
\item {\bf Photo Editing:} Auto-enhance, remove backgrounds
\item {\bf Gaming:} Opponents that adapt to your skill level
\item {\bf Movie Recommendations:} Netflix knows what you'll enjoy
\item {\bf AI-Generated Art:} Create images from text descriptions
\end{itemize}

Example: Type ``sunset over mountains with purple clouds'' and AI creates it!
\end{frame}

%%%%%%%%%%%%%%%%%%%%%%%%%%%%%%%%%%%%%%%%%%%%%%%%%%%%%%%%%%%%%%%%%%%%%%%%%%%%%%%%%%
\begin{frame}[fragile]\frametitle{AI in Agriculture}

\begin{itemize}
\item {\bf Crop Monitoring:} Drones check plant health from above
\item {\bf Weather Prediction:} Better forecasts for planting
\item {\bf Pest Detection:} Identifies diseases early
\item {\bf Irrigation:} Waters only when needed, saves water
\item {\bf Harvest Timing:} Predicts best time to harvest
\end{itemize}

Helping farmers increase yields and reduce waste!
\end{frame}

%%%%%%%%%%%%%%%%%%%%%%%%%%%%%%%%%%%%%%%%%%%%%%%%%%%%%%%%%%%%%%%%%%%%%%%%%%%%%%%%%%
\begin{frame}[fragile]\frametitle{}
\begin{center}
{\Large Common Questions and Concerns}
\end{center}
\end{frame}

%%%%%%%%%%%%%%%%%%%%%%%%%%%%%%%%%%%%%%%%%%%%%%%%%%%%%%%%%%%%%%%%%%%%%%%%%%%%%%%%%%
\begin{frame}[fragile]\frametitle{``Will AI Take My Job?''}

{\bf Honest Answer:} Some jobs will change, some will disappear

{\bf Jobs Most Affected:}
\begin{itemize}
\item Repetitive tasks: Data entry, assembly line work
\item Some customer service (chatbots)
\item Basic bookkeeping
\end{itemize}

{\bf Jobs Safe (and Growing):}
\begin{itemize}
\item Anything requiring human touch: nursing, teaching, counseling
\item Creative work: design, art, writing
\item Complex decision-making: management, strategy
\item Working WITH AI: new tech roles
\end{itemize}

{\bf Key:} Learn to work alongside AI, not compete with it
\end{frame}

%%%%%%%%%%%%%%%%%%%%%%%%%%%%%%%%%%%%%%%%%%%%%%%%%%%%%%%%%%%%%%%%%%%%%%%%%%%%%%%%%%
\begin{frame}[fragile]\frametitle{Real Examples: Jobs Transformed}

{\bf Bank Teller:}
\begin{itemize}
\item Before: Processed transactions all day
\item Now: ATMs do basic transactions, tellers help with complex needs
\item Better job: More customer interaction, less repetitive work
\end{itemize}

{\bf Travel Agent:}
\begin{itemize}
\item Before: Book flights, hotels
\item Now: Websites do simple bookings, agents plan complex trips
\item Better job: Expertise-based, more interesting
\end{itemize}

Pattern: AI does boring parts, humans do interesting parts!
\end{frame}

%%%%%%%%%%%%%%%%%%%%%%%%%%%%%%%%%%%%%%%%%%%%%%%%%%%%%%%%%%%%%%%%%%%%%%%%%%%%%%%%%%
\begin{frame}[fragile]\frametitle{``Is My Personal Data Safe?''}

{\bf The Reality:}
\begin{itemize}
\item AI needs data to learn and improve
\item Your data = photos, messages, search history, purchases
\item Companies use this to make better services
\end{itemize}

{\bf What You Can Do:}
\begin{itemize}
\item Read privacy settings (yes, actually read them!)
\item Limit what you share on social media
\item Use privacy-focused browsers
\item Turn off location when not needed
\item Regularly delete search history
\end{itemize}

Balance: Convenience vs Privacy - YOU decide!
\end{frame}

%%%%%%%%%%%%%%%%%%%%%%%%%%%%%%%%%%%%%%%%%%%%%%%%%%%%%%%%%%%%%%%%%%%%%%%%%%%%%%%%%%
\begin{frame}[fragile]\frametitle{``Can AI Make Mistakes?''}

{\bf Yes! AI makes mistakes, sometimes serious ones:}

\begin{itemize}
\item Face recognition fails for certain skin tones
\item Language translators make funny (or offensive) errors
\item Self-driving cars can't handle all situations
\item Medical AI might miss rare conditions
\item Loan approval AI might be biased
\end{itemize}

{\bf Remember:}
\begin{itemize}
\item AI learns from data
\item If data has bias, AI will have bias
\item Always have humans check important decisions
\end{itemize}
\end{frame}

%%%%%%%%%%%%%%%%%%%%%%%%%%%%%%%%%%%%%%%%%%%%%%%%%%%%%%%%%%%%%%%%%%%%%%%%%%%%%%%%%%
\begin{frame}[fragile]\frametitle{``Will Robots Become Like Humans?''}

{\bf What you see in movies:} Robots with feelings, consciousness

{\bf Current Reality:}
\begin{itemize}
\item AI can seem intelligent but isn't ``conscious''
\item Doesn't have feelings, desires, or awareness
\item Can't truly understand context like humans
\item Limited to what it's trained for
\end{itemize}

{\bf Future:}
\begin{itemize}
\item Scientists debate if this is even possible
\item If it happens, it's decades away
\item For now, AI is a tool, not a thinking being
\end{itemize}
\end{frame}

%%%%%%%%%%%%%%%%%%%%%%%%%%%%%%%%%%%%%%%%%%%%%%%%%%%%%%%%%%%%%%%%%%%%%%%%%%%%%%%%%%
\begin{frame}[fragile]\frametitle{}
\begin{center}
{\Large How to Use AI Wisely in Your Life}
\end{center}
\end{frame}

%%%%%%%%%%%%%%%%%%%%%%%%%%%%%%%%%%%%%%%%%%%%%%%%%%%%%%%%%%%%%%%%%%%%%%%%%%%%%%%%%%
\begin{frame}[fragile]\frametitle{Practical Tips for Everyone}

{\bf DO:}
\begin{itemize}
\item Use AI to save time on routine tasks
\item Let AI handle boring, repetitive work
\item Ask AI to explain things you don't understand
\item Use AI recommendations as suggestions, not commands
\item Verify important information from multiple sources
\end{itemize}

{\bf DON'T:}
\begin{itemize}
\item Blindly trust everything AI says
\item Share overly personal information
\item Rely completely on AI for critical decisions
\item Use AI for unethical purposes
\end{itemize}
\end{frame}

%%%%%%%%%%%%%%%%%%%%%%%%%%%%%%%%%%%%%%%%%%%%%%%%%%%%%%%%%%%%%%%%%%%%%%%%%%%%%%%%%%
\begin{frame}[fragile]\frametitle{Everyday Uses That Make Life Easier}

{\bf For Home Management:}
\begin{itemize}
\item Ask AI for meal plans based on your preferences
\item Get cleaning schedules and tips
\item Budget tracking and expense suggestions
\end{itemize}

{\bf For Shopping:}
\begin{itemize}
\item Compare prices across websites
\item Get gift suggestions
\item Find best deals and discounts
\end{itemize}

{\bf For Learning:}
\begin{itemize}
\item Learn new recipes, crafts, skills
\item Get explanations in simple language
\item Practice new languages
\end{itemize}
\end{frame}

%%%%%%%%%%%%%%%%%%%%%%%%%%%%%%%%%%%%%%%%%%%%%%%%%%%%%%%%%%%%%%%%%%%%%%%%%%%%%%%%%%
\begin{frame}[fragile]\frametitle{Safety Tips for Your Family}

{\bf For Children:}
\begin{itemize}
\item Monitor their AI interactions
\item Teach them not to share personal info
\item Explain that AI isn't always right
\item Encourage critical thinking
\end{itemize}

{\bf For Elderly:}
\begin{itemize}
\item Voice assistants can help with reminders
\item But watch out for scams using AI voices
\item Always verify before clicking links
\item Set up trusted contacts for verification
\end{itemize}
\end{frame}

%%%%%%%%%%%%%%%%%%%%%%%%%%%%%%%%%%%%%%%%%%%%%%%%%%%%%%%%%%%%%%%%%%%%%%%%%%%%%%%%%%
\begin{frame}[fragile]\frametitle{}
\begin{center}
{\Large The Positive Side - How AI Helps Society}
\end{center}
\end{frame}

%%%%%%%%%%%%%%%%%%%%%%%%%%%%%%%%%%%%%%%%%%%%%%%%%%%%%%%%%%%%%%%%%%%%%%%%%%%%%%%%%%
\begin{frame}[fragile]\frametitle{AI Making the World Better}

{\bf Environment:}
\begin{itemize}
\item Predicting natural disasters
\item Optimizing energy usage in cities
\item Tracking endangered animals
\item Reducing food waste in supply chains
\end{itemize}

{\bf Social Good:}
\begin{itemize}
\item Making education accessible to remote areas
\item Translating for refugees
\item Helping disabled people communicate
\item Detecting child exploitation online
\end{itemize}
\end{frame}

%%%%%%%%%%%%%%%%%%%%%%%%%%%%%%%%%%%%%%%%%%%%%%%%%%%%%%%%%%%%%%%%%%%%%%%%%%%%%%%%%%
\begin{frame}[fragile]\frametitle{Inspiring Real Stories}

{\bf Story 1: Finding Missing Children}
\begin{itemize}
\item AI analyzes thousands of photos instantly
\item Helps find missing children faster
\end{itemize}

{\bf Story 2: Helping Farmers}
\begin{itemize}
\item Indian farmers use AI apps in local languages
\item Get crop advice, weather alerts on basic phones
\item Increased yields by 30\%
\end{itemize}

{\bf Story 3: Breaking Language Barriers}
\begin{itemize}
\item Real-time translation helps doctors treat foreign patients
\item Breaking down communication barriers worldwide
\end{itemize}
\end{frame}

%%%%%%%%%%%%%%%%%%%%%%%%%%%%%%%%%%%%%%%%%%%%%%%%%%%%%%%%%%%%%%%%%%%%%%%%%%%%%%%%%%
\begin{frame}[fragile]\frametitle{}
\begin{center}
{\Large Getting Comfortable with AI}
\end{center}
\end{frame}

%%%%%%%%%%%%%%%%%%%%%%%%%%%%%%%%%%%%%%%%%%%%%%%%%%%%%%%%%%%%%%%%%%%%%%%%%%%%%%%%%%
\begin{frame}[fragile]\frametitle{Start Small - Try These Today}

{\bf Beginner Level:}
\begin{itemize}
\item Use Google Assistant/Siri for simple tasks
\item Let Gmail organize your inbox
\item Try Netflix/YouTube recommendations
\end{itemize}

{\bf Intermediate Level:}
\begin{itemize}
\item Chat with ChatGPT (it's free!)
\item Use AI photo editors
\item Try Google Lens (point camera, get info)
\end{itemize}

{\bf Advanced:}
\begin{itemize}
\item Use AI for work presentations
\item Explore AI art generation
\item Learn basic prompt writing (how to ask AI effectively)
\end{itemize}
\end{frame}

%%%%%%%%%%%%%%%%%%%%%%%%%%%%%%%%%%%%%%%%%%%%%%%%%%%%%%%%%%%%%%%%%%%%%%%%%%%%%%%%%%
\begin{frame}[fragile]\frametitle{Free AI Tools You Can Try Right Now}

\begin{itemize}
\item {\bf ChatGPT:} Conversation and assistance
\item {\bf Google Bard:} Similar to ChatGPT
\item {\bf DALL-E Mini:} Create images from text
\item {\bf Grammarly:} Writing improvement
\item {\bf Google Translate:} Better than ever with AI
\item {\bf Canva Magic Write:} Design help
\item {\bf Microsoft Bing Chat:} Search with conversation
\end{itemize}

All available on your phone or computer!
\end{frame}

%%%%%%%%%%%%%%%%%%%%%%%%%%%%%%%%%%%%%%%%%%%%%%%%%%%%%%%%%%%%%%%%%%%%%%%%%%%%%%%%%%
% [ai-for-x: commented 2026-11-07 | covered by ai_core_nontech_prompts.tex]
% \begin{frame}[fragile]\frametitle{How to Talk to AI (Prompting)}
%
% {\bf Poor Prompt:} ``Tell me about food''
% \begin{itemize}
% \item Too vague, you'll get generic answer
% \end{itemize}
%
% {\bf Good Prompt:} ``Suggest 3 healthy dinner recipes using chicken, vegetables, and Indian spices that can be cooked in 30 minutes''
% \begin{itemize}
% \item Specific, clear, actionable
% \end{itemize}
%
% {\bf Tips:}
% \begin{itemize}
% \item Be specific about what you want
% \item Provide context
% \item Ask for format you prefer (list, paragraph, steps)
% \item Can ask AI to explain like you're 5 years old!
% \end{itemize}
% \end{frame}

%%%%%%%%%%%%%%%%%%%%%%%%%%%%%%%%%%%%%%%%%%%%%%%%%%%%%%%%%%%%%%%%%%%%%%%%%%%%%%%%%%
\begin{frame}[fragile]\frametitle{}
\begin{center}
{\Large What's Next? The Future}
\end{center}
\end{frame}

%%%%%%%%%%%%%%%%%%%%%%%%%%%%%%%%%%%%%%%%%%%%%%%%%%%%%%%%%%%%%%%%%%%%%%%%%%%%%%%%%%
\begin{frame}[fragile]\frametitle{What to Expect in Next 5-10 Years}

{\bf At Home:}
\begin{itemize}
\item Fully automated smart homes
\item AI companions for elderly care
\item Personalized health monitoring
\end{itemize}

{\bf In Cities:}
\begin{itemize}
\item Self-driving taxis common
\item Traffic lights that optimize flow
\item Energy grids that balance themselves
\end{itemize}

{\bf At Work:}
\begin{itemize}
\item AI assistants for everyone
\item More remote work enabled by AI
\item Focus on creative, strategic work
\end{itemize}
\end{frame}

%%%%%%%%%%%%%%%%%%%%%%%%%%%%%%%%%%%%%%%%%%%%%%%%%%%%%%%%%%%%%%%%%%%%%%%%%%%%%%%%%%
\begin{frame}[fragile]\frametitle{Preparing Yourself and Family}

{\bf Stay Informed:}
\begin{itemize}
\item Follow tech news in simple language
\item Ask younger family members to teach you
\item Join community classes
\end{itemize}

{\bf Be Adaptable:}
\begin{itemize}
\item Don't fear change
\item Try new things gradually
\item Learn basic digital skills
\end{itemize}

{\bf Stay Human:}
\begin{itemize}
\item Maintain personal connections
\item Don't let AI replace human interaction
\item Use technology to enhance life, not replace it
\end{itemize}
\end{frame}

%%%%%%%%%%%%%%%%%%%%%%%%%%%%%%%%%%%%%%%%%%%%%%%%%%%%%%%%%%%%%%%%%%%%%%%%%%%%%%%%%%
\begin{frame}[fragile]\frametitle{}
\begin{center}
{\Large Key Takeaways}
\end{center}
\end{frame}

%%%%%%%%%%%%%%%%%%%%%%%%%%%%%%%%%%%%%%%%%%%%%%%%%%%%%%%%%%%%%%%%%%%%%%%%%%%%%%%%%%
\begin{frame}[fragile]\frametitle{Remember These Important Points}

\begin{enumerate}
\item AI is already helping you daily (often invisibly)
\item AI learns from examples, like humans do
\item It's a powerful tool, but not magic
\item Use it to save time, not replace thinking
\item Your job may change, but you can adapt
\item Privacy matters - be conscious about data
\item AI makes mistakes - always verify important info
\item The future is AI-assisted, not AI-replaced
\item You can learn to use AI at any age
\item Human connection remains irreplaceable
\end{enumerate}

\end{frame}

%%%%%%%%%%%%%%%%%%%%%%%%%%%%%%%%%%%%%%%%%%%%%%%%%%%%%%%%%%%%%%%%%%%%%%%%%%%%%%%%%%
\begin{frame}[fragile]\frametitle{Final Thoughts}

\begin{center}
\Large

AI is not here to replace you

\vspace{0.5cm}

It's here to {\bf work with you}

\vspace{0.5cm}

The best future is one where {\bf humans and AI collaborate}

\vspace{0.5cm}

You bring: creativity, empathy, wisdom, judgment

AI brings: speed, memory, pattern recognition

\vspace{0.5cm}

{\bf Together = Powerful!}

\end{center}

\end{frame}

%%%%%%%%%%%%%%%%%%%%%%%%%%%%%%%%%%%%%%%%%%%%%%%%%%%%%%%%%%%%%%%%%%%%%%%%%%%%%%%%%%
\begin{frame}[fragile]\frametitle{Your Next Steps}

{\bf This Week:}
\begin{itemize}
\item Try one AI tool (ChatGPT or voice assistant)
\item Notice where AI is already in your life
\end{itemize}

{\bf This Month:}
\begin{itemize}
\item Use AI for one practical task
\item Share what you learned with family
\end{itemize}

{\bf This Year:}
\begin{itemize}
\item Become comfortable with AI basics
\item Help others in your community understand AI
\end{itemize}

\vspace{0.5cm}
Remember: It's never too late to start!
\end{frame}

%%%%%%%%%%%%%%%%%%%%%%%%%%%%%%%%%%%%%%%%%%%%%%%%%%%%%%%%%%%%%%%%%%%%%%%%%%%%%%%%%%
\begin{frame}[fragile]\frametitle{Questions to Discuss}

Think about and discuss with family/friends:

\begin{itemize}
\item Where have you unknowingly used AI today?
\item What's one task you'd like AI to help you with?
\item What worries you most about AI?
\item How can AI make your community better?
\item What's one thing AI should never do?
\end{itemize}

\vspace{0.5cm}
These conversations help us shape how AI develops!
\end{frame}

% %%%%%%%%%%%%%%%%%%%%%%%%%%%%%%%%%%%%%%%%%%%%%%%%%%%%%%%%%%%%%%%%%%%%%%%%%%%%%%%%%%
% \begin{frame}[fragile]\frametitle{}
% \begin{center}
% {\Huge Thank You!}

% \vspace{1cm}
% {\Large Embrace the Future with Confidence}

% \vspace{0.5cm}
% Remember: Technology should serve humanity, not the other way around.

% \vspace{0.5cm}
% Stay curious, stay informed, stay human!
% \end{center}
% \end{frame}

% %%%%%%%%%%%%%%%%%%%%%%%%%%%%%%%%%%%%%%%%%%%%%%%%%%%%%%%%%%%%%%%%%%%%%%%%%%%%%%%%%%
% \begin{frame}[fragile]\frametitle{Quick Reference Summary}

% {\bf What is AI?} Teaching machines to be smart

% {\bf Where is it?} Gmail, Netflix, Google Maps, Banking, Shopping

% {\bf How does it learn?} From examples and practice, like humans

% {\bf Is it dangerous?} Can be, if misused. Use responsibly

% {\bf Will it take jobs?} Some will change, new ones will come

% {\bf Should I learn it?} Yes! Start with simple tools

% {\bf Can I trust it?} Verify important information always

% {\bf What's the future?} AI-assisted everything, but humans remain essential

% \end{frame}

\section[fluency]{Fluency}
%%%%%%%%%%%%%%%%%%%%%%%%%%%%%%%%%%%%%%%%%%%%%%%%%%%%%%%%%%%%%%%%%%%%%%%%%%%%%%%%%%
\begin{frame}[fragile]\frametitle{}
\begin{center}
{\Large AI Fluency Framework}

{\tiny (Ref: Anthropic Academy Course)}
\end{center}
\end{frame}

%%%%%%%%%%%%%%%%%%%%%%%%%%%%%%%%%%%%%%%%%%%%%%%%%%%%%%%%%%%%%%%%%%%%%%%%%%%%%%%%%%
\begin{frame}[fragile]\frametitle{What is AI Fluency?}
\begin{columns}
    \begin{column}[T]{0.6\linewidth}
      \begin{itemize}
        \item \textbf{AI Fluency} is the ability to work with AI systems in ways that are:
          \begin{itemize}
            \item Effective: achieving your intended goals
            \item Efficient: making good use of time and resources
            \item Ethical: respecting values and responsibilities
            \item Safe: avoiding harm to people and systems
          \end{itemize}
        \item It combines:
          \begin{itemize}
            \item Practical skills and domain knowledge
            \item Critical insight and contextual judgment
            \item Values that guide responsible AI use
            \item Adaptability to evolving AI technologies
          \end{itemize}
      \end{itemize}
    \end{column}
    \begin{column}[T]{0.4\linewidth}
      \begin{block}{The 4Ds Framework}
        The four core competencies of AI Fluency:
        \begin{enumerate}
          \item \textbf{Delegation}
          \item \textbf{Description}
          \item \textbf{Discernment}
          \item \textbf{Diligence}
        \end{enumerate}
      \end{block}
    \end{column}
  \end{columns}
\end{frame}

%%%%%%%%%%%%%%%%%%%%%%%%%%%%%%%%%%%%%%%%%%%%%%%%%%%%%%%%%%%%%%%%%%%%%%%%%%%%%%%%%%
\begin{frame}[fragile]\frametitle{The 4Ds: Delegation}
\begin{columns}
    \begin{column}[T]{0.6\linewidth}
      \begin{itemize}
        \item \textbf{Delegation}: Deciding what work should be done by humans, by AI, or by both together
        \item Key sub-skills:
          \begin{itemize}
            \item \textbf{Problem Awareness:} Clearly understanding your goals and the nature of the work \emph{before} involving AI
            \item \textbf{Platform Awareness:} Understanding the capabilities and limitations of different AI systems
            \item \textbf{Task Delegation:} Thoughtfully distributing work between humans and AI to leverage the strengths of each
          \end{itemize}
        \item Good delegation means neither under-using nor over-relying on AI
      \end{itemize}
    \end{column}
    \begin{column}[T]{0.4\linewidth}
      \begin{alertblock}{Key Question}
        \emph{``Which parts of this task are best handled by AI, and which require human judgment?''}
      \end{alertblock}
      \vspace{0.5em}
      \begin{exampleblock}{Example}
        Delegating first-draft writing to AI while retaining human oversight for factual accuracy and tone
      \end{exampleblock}
    \end{column}
  \end{columns}
\end{frame}

%%%%%%%%%%%%%%%%%%%%%%%%%%%%%%%%%%%%%%%%%%%%%%%%%%%%%%%%%%%%%%%%%%%%%%%%%%%%%%%%%%
\begin{frame}[fragile]\frametitle{The 4Ds: Description}
\begin{columns}
    \begin{column}[T]{0.6\linewidth}
      \begin{itemize}
        \item \textbf{Description}: Effectively communicating with AI systems
        \item Three dimensions:
          \begin{itemize}
            \item \textbf{Product Description:} Define the desired \emph{output}: format, audience, style, length
            \item \textbf{Process Description:} Define \emph{how} the AI should approach the task: step-by-step instructions, constraints, order of operations
            \item \textbf{Performance Description:} Define the AI's \emph{behaviour} during the interaction: concise vs. detailed, challenging vs. supportive
          \end{itemize}
        \item Clear description reduces guesswork and improves output quality
      \end{itemize}
    \end{column}
    \begin{column}[T]{0.4\linewidth}
      \begin{alertblock}{Key Question}
        \emph{``Have I given the AI enough context to know what I actually need?''}
      \end{alertblock}
      \vspace{0.5em}
      \begin{exampleblock}{Example}
        ``Summarise this report in three bullet points for a non-technical executive audience, using plain language''
      \end{exampleblock}
    \end{column}
  \end{columns}
\end{frame}

%%%%%%%%%%%%%%%%%%%%%%%%%%%%%%%%%%%%%%%%%%%%%%%%%%%%%%%%%%%%%%%%%%%%%%%%%%%%%%%%%%
\begin{frame}[fragile]\frametitle{The 4Ds: Discernment}
\begin{columns}
    \begin{column}[T]{0.6\linewidth}
      \begin{itemize}
        \item \textbf{Discernment}: Thoughtfully and critically evaluating AI outputs, processes, and behaviour
        \item Three dimensions:
          \begin{itemize}
            \item \textbf{Product Discernment:} Evaluate the \emph{quality} of what AI produces: accuracy, appropriateness, coherence, relevance
            \item \textbf{Process Discernment:} Evaluate \emph{how} the AI arrived at its output: check for logical errors, lapses in attention, or flawed reasoning steps
            \item \textbf{Performance Discernment:} Evaluate \emph{how} the AI behaves during interaction: is its communication style effective for your needs?
          \end{itemize}
      \end{itemize}
    \end{column}
    \begin{column}[T]{0.4\linewidth}
      \begin{alertblock}{Key Question}
        \emph{``Is this output accurate, appropriate, and free from subtle errors?''}
      \end{alertblock}
      \vspace{0.5em}
      \begin{exampleblock}{Watch out for}
        \begin{itemize}
          \item Hallucinations
          \item Plausible-sounding but incorrect facts
          \item Bias in generated content
        \end{itemize}
      \end{exampleblock}
    \end{column}
  \end{columns}
\end{frame}

%%%%%%%%%%%%%%%%%%%%%%%%%%%%%%%%%%%%%%%%%%%%%%%%%%%%%%%%%%%%%%%%%%%%%%%%%%%%%%%%%%
\begin{frame}[fragile]\frametitle{The 4Ds: Diligence}
\begin{columns}
    \begin{column}[T]{0.6\linewidth}
      \begin{itemize}
        \item \textbf{Diligence}: Using AI responsibly and ethically
        \item Three dimensions:
          \begin{itemize}
            \item \textbf{Creation Diligence:} Be thoughtful about \emph{which} AI systems you use and how you interact with them
            \item \textbf{Transparency Diligence:} Be honest about AI's role in your work with everyone who needs to know
            \item \textbf{Deployment Diligence:} Take responsibility for verifying and vouching for the outputs you use or share
          \end{itemize}
        \item Diligence means you remain \emph{accountable} for AI-assisted work: the AI is a tool, not the decision-maker
      \end{itemize}
    \end{column}
    \begin{column}[T]{0.4\linewidth}
      \begin{alertblock}{Key Question}
        \emph{``Am I using AI in a way I would be comfortable disclosing and defending?''}
      \end{alertblock}
      \vspace{0.5em}
      \begin{exampleblock}{Example}
        Disclosing AI assistance in a report, and verifying all facts before publication
      \end{exampleblock}
    \end{column}
  \end{columns}
\end{frame}

%%%%%%%%%%%%%%%%%%%%%%%%%%%%%%%%%%%%%%%%%%%%%%%%%%%%%%%%%%%%%%%%%%%%%%%%%%%%%%%%%%
\begin{frame}[fragile]\frametitle{Human--AI Interaction Modes}
\begin{columns}
    \begin{column}[T]{0.55\linewidth}
      \begin{itemize}
        \item \textbf{Automation}
          \begin{itemize}
            \item AI performs specific tasks based on specific human instructions
            \item Human defines \emph{what} needs to be done; AI executes it
            \item Example: generating a formatted report from structured data
          \end{itemize}
        \item \textbf{Augmentation}
          \begin{itemize}
            \item Humans and AI collaborate as thinking partners
            \item Iterative back-and-forth; both contribute to the outcome
            \item Example: co-writing and refining a document together
          \end{itemize}
        \item \textbf{Agency}
          \begin{itemize}
            \item AI works independently on the human's behalf
            \item Human sets goals and behaviour patterns, not exact actions
            \item Example: an AI agent that schedules meetings autonomously
          \end{itemize}
      \end{itemize}
    \end{column}
    \begin{column}[T]{0.45\linewidth}
      \begin{block}{Spectrum of Control}
        \begin{center}
          \begin{tabular}{cl}
            \textbf{Mode} & \textbf{Human role} \\
            \hline
            Automation  & Director \\
            Augmentation & Collaborator \\
            Agency      & Goal-setter \\
          \end{tabular}
        \end{center}
      \end{block}
      \vspace{0.4em}
      Choosing the right mode is itself a \textbf{Delegation} skill.
    \end{column}
  \end{columns}
\end{frame}

%%%%%%%%%%%%%%%%%%%%%%%%%%%%%%%%%%%%%%%%%%%%%%%%%%%%%%%%%%%%%%%%%%%%%%%%%%%%%%%%%%
\begin{frame}[fragile]\frametitle{AI Technical Concepts (1/2)}
% \begin{columns}
    % \begin{column}[T]{0.5\linewidth}
      \begin{itemize}
        \item \textbf{Generative AI}: AI systems that \emph{create} new content (text, images, code) rather than merely analysing existing data
        \item \textbf{Large Language Models (LLMs)}: Generative AI trained on vast text corpora to understand and generate human language; e.g.\ Claude
        \item \textbf{Parameters}: Mathematical values within a model that determine how it processes information; modern LLMs contain \emph{billions}
        \item \textbf{Neural Networks}: Interconnected nodes in layers that learn patterns from data; inspired by, but distinct from, biological brains
        \item \textbf{Transformer Architecture}: Breakthrough 2017 design enabling LLMs to process text in parallel while tracking relationships across long passages
      % \end{itemize}
    % \end{column}
    % \begin{column}[T]{0.5\linewidth}
      % \begin{itemize}
        \item \textbf{Pre-training}: Initial phase where a model learns patterns from massive text datasets, building foundational language understanding
        \item \textbf{Fine-tuning}: Additional training to make a model follow instructions, give helpful responses, and avoid harmful content
        \item \textbf{Scaling Laws}: Empirical observation that as models grow larger and train on more data, performance improves consistently; new capabilities can \emph{emerge} at scale thresholds
        \item \textbf{Reasoning / Thinking Models}: Models designed to work step-by-step through complex problems, showing improved logical reasoning
      \end{itemize}
    % \end{column}
  % \end{columns}
\end{frame}

%%%%%%%%%%%%%%%%%%%%%%%%%%%%%%%%%%%%%%%%%%%%%%%%%%%%%%%%%%%%%%%%%%%%%%%%%%%%%%%%%%
\begin{frame}[fragile]\frametitle{AI Technical Concepts (2/2)}
% \begin{columns}
    % \begin{column}[T]{0.5\linewidth}
      \begin{itemize}
        \item \textbf{Context Window}: The amount of information an AI can consider at one time (conversation history + documents); has a maximum limit that varies by model
        \item \textbf{Hallucination}: An error where the AI confidently states something that sounds plausible but is actually incorrect
        \item \textbf{Knowledge Cutoff Date}: The point after which the model has no built-in knowledge, based on when training data ended
        \item \textbf{Temperature}: A setting controlling output randomness: higher temperature $\rightarrow$ more creative; lower temperature $\rightarrow$ more predictable
      % \end{itemize}
    % \end{column}
    % \begin{column}[T]{0.5\linewidth}
      % \begin{itemize}
        \item \textbf{Retrieval Augmented Generation (RAG)}: Connects AI to external knowledge sources at inference time, improving accuracy and reducing hallucinations
        \item \textbf{Bias}: Systematic patterns in AI outputs that unfairly favour or disadvantage certain groups, often reflecting imbalances in training data
      \end{itemize}
      \vspace{0.4em}
      \begin{alertblock}{Practical implication}
        Hallucination + knowledge cutoff + bias together mean \textbf{Discernment} is always necessary: never accept AI output uncritically.
      \end{alertblock}
    % \end{column}
  % \end{columns}
\end{frame}

%%%%%%%%%%%%%%%%%%%%%%%%%%%%%%%%%%%%%%%%%%%%%%%%%%%%%%%%%%%%%%%%%%%%%%%%%%%%%%%%%%
% [ai-for-x: commented 2026-11-07 | covered by ai_core_nontech_prompts.tex]
% \begin{frame}[fragile]\frametitle{Prompt Engineering Concepts}
% % \begin{columns}
%     % \begin{column}[T]{0.5\linewidth}
%       \begin{itemize}
%         \item \textbf{Prompt}: The full input given to an AI model: instructions, context, documents, and examples
%         \item \textbf{Prompt Engineering}: The practice of designing effective prompts to produce desired outputs; combines clear communication with AI-specific techniques
%         \item \textbf{Chain-of-Thought Prompting}: Encouraging the AI to work step-by-step through a problem before giving a final answer; improves accuracy on complex tasks
%         \item \textbf{Few-Shot Learning (n-shot prompting)}: Providing $n$ input--output examples so the model understands the desired pattern without lengthy explanations
%       % \end{itemize}
%     % \end{column}
%     % \begin{column}[T]{0.5\linewidth}
%       % \begin{itemize}
%         \item \textbf{Role / Persona Definition}: Specifying a character, expertise level, or communication style for the AI to adopt (e.g.\ ``Explain this as a UX expert'')
%         \item \textbf{Output Constraints / Formatting}: Specifying the desired format, length, or structure within the prompt% (e.g.\ ``Reply in a 3-column table, max 200 words'')
%         \item \textbf{Think-First Approach}: Explicitly asking the AI to reason through its thinking \emph{before} providing a final answer, leading to more considered responses
%       \end{itemize}
%     % \end{column}
%   % \end{columns}
% \end{frame}

%%%%%%%%%%%%%%%%%%%%%%%%%%%%%%%%%%%%%%%%%%%%%%%%%%%%%%%%%%%%%%%%%%%%%%%%%%%%%%%%%%
\begin{frame}[fragile]\frametitle{Putting It All Together}
\begin{columns}
    \begin{column}[T]{0.6\linewidth}
      \begin{itemize}
        \item AI Fluency is not a single skill: it is a \textbf{framework} that combines technical understanding, communication, critical thinking, and ethical responsibility
        \item The \textbf{4Ds} operate as a cycle:
          \begin{enumerate}
            \item \textit{Delegate} the right tasks to AI
            \item \textit{Describe} your needs clearly
            \item \textit{Discern} the quality of outputs
            \item \textit{Diligence}: take responsibility for what you share
          \end{enumerate}
        \item Understanding technical concepts (LLMs, hallucination, context windows) helps you apply the 4Ds more effectively
        \item Prompt engineering skills make your \textbf{Description} more powerful
      \end{itemize}
    \end{column}
    \begin{column}[T]{0.4\linewidth}
      \begin{block}{Remember}
        \begin{itemize}
          \item AI is a \emph{tool}, not an authority
          \item You remain \emph{accountable} for AI-assisted outputs
          \item Fluency grows with \emph{practice} and reflection
        \end{itemize}
      \end{block}
      \vspace{0.5em}
      {\tiny Source: Rick Dakan, Joseph Feller \& Anthropic, \textit{AI Fluency: Key Terminology Cheat Sheet}, 2025. CC BY-NC-SA 4.0}
    \end{column}
  \end{columns}
\end{frame}

%%%%%%%%%%%%%%%%%%%%%%%%%%%%%%%%%%%%%%%%%%%%%%%%%%%%%%%%%%%%%%%%%%%%%%%%%%%%%%%%%%

% ChatGPT section folded in from the retired Main_Seminar_AI_ChatGPT_NonTech standalone driver
% (Aug 2026 consolidation). Deliberately keeps chatgpt_intro_short (the simplified telling), not
% the full chatgpt_intro used in AI_For_All_Tech -- appropriate for this audience.
\section[ChatGPT]{ChatGPT}
%%%%%%%%%%%%%%%%%%%%%%%%%%%%%%%%%%%%%%%%%%%%%%%%%%%%%%%%%%%%%%%%%%%%%%%%%%%%%%%%%%
\begin{frame}[fragile]\frametitle{}
\begin{center}
{\Large Overview of ChatGPT}
\end{center}
\end{frame}

%%%%%%%%%%%%%%%%%%%%%%%%%%%%%%%%%%%%%%%%%%%%%%%%%%%%%%%%%%%
\begin{frame}[fragile]\frametitle{What is ChatGPT?}


\begin{itemize}
\item A Chatbot
\item A language model, that takes 'prompt' from user and generates a response.
\item Built by OpenAI and released in Nov 2022
\item Got 1m users in 5 days  (Insta 2.5 months, Spotify 5m, Facebook 10m, Netflix 3.5yrs)
\end{itemize}	 

\begin{itemize}
\item Link: https://chat.openai.com/chat
\item Details: https://openai.com/blog/chatgpt/
\end{itemize}	 

\begin{center}
\includegraphics[width=0.5\linewidth,keepaspectratio]{chatgpt33}
\end{center}		
		
{\tiny (Ref: OpenAI, creator of ChatGPT, casts spell on Microsoft
- The Jakarta Post}
			
			
\end{frame}

% %%%%%%%%%%%%%%%%%%%%%%%%%%%%%%%%%%%%%%%%%%%%%%%%%%%%%%%%%%%
% \begin{frame}[fragile]\frametitle{What is ChatGPT?}

% ChatGPT is powered by a Large Language Model (GPT-3.5) optimized for 
% dialogue, developed by OpenAI

% \begin{itemize}
% \item Link: https://chat.openai.com/chat
% \item Details: https://openai.com/blog/chatgpt/
% \end{itemize}	 

% \begin{center}
% \includegraphics[width=0.6\linewidth,keepaspectratio]{chatgpt33}
% \end{center}		
		
% {\tiny (Ref: OpenAI, creator of ChatGPT, casts spell on Microsoft
% - The Jakarta Post}

% \end{frame}

%%%%%%%%%%%%%%%%%%%%%%%%%%%%%%%%%%%%%%%%%%%%%%%%%%%%%%%%%%%
% [ai-for-x: commented 2026-11-07 | covered by ai_core_nontech_whatisai.tex]
% \begin{frame}[fragile]\frametitle{What is a Language Models?}
%
% \begin{itemize}
% \item While typing SMS, have you seen it suggests next word?
% \item While typing email, have you seen next few words are suggested?
% \item How does it suggest? (suggestions are not random, right?)
% \item In the past, for ``Lets go for a \ldots', if you have typed 'coffee' 15 times, 'movie' say 4 times, then it learns that. Machine/Statistical Learning.
% \item Next time, when you type ``Lets go for a '', what will be suggested? why?
% \item This is called Language Model. Predicting the next word. When done continuously, one after other, it spits sentence, called Generative Model.
% \end{itemize}
%
% \begin{center}
% \includegraphics[width=0.6\linewidth,keepaspectratio]{chatgpt34}
% \end{center}
%
% \end{frame}

% %%%%%%%%%%%%%%%%%%%%%%%%%%%%%%%%%%%%%%%%%%%%%%%%%%%%%%%%%%%
% \begin{frame}[fragile]\frametitle{Evolution of Language Models}

% Language Models can be statistical (frequency based) or Machine/Deep Learning (supervised) based. Simple to complex.

% \begin{center}
% \includegraphics[width=\linewidth,keepaspectratio]{chatgpt30}
% \end{center}				
% {\tiny (Ref: Analytics Vidhya https://editor.analyticsvidhya.com/uploads/59483evolution\_of\_NLP.png)}

% \end{frame}

%%%%%%%%%%%%%%%%%%%%%%%%%%%%%%%%%%%%%%%%%%%%%%%%%%%%%%%%%%%
% [ai-for-x: commented 2026-11-07 | covered by ai_core_nontech_whatisai.tex]
% \begin{frame}[fragile]\frametitle{Large Language Models - Comparison}
%
% \begin{center}
% \includegraphics[width=\linewidth,keepaspectratio]{chatgpt31}
% \end{center}
% {\tiny (Ref: Deus.ai https://www.deus.ai/post/gpt-3-what-is-all-the-excitement-about)}
%
% \end{frame}



% %%%%%%%%%%%%%%%%%%%%%%%%%%%%%%%%%%%%%%%%%%%%%%%%%%%%%%%%%%%
% \begin{frame}[fragile]\frametitle{What is ChatGPT?}


% \begin{itemize}
% \item A Chatbot
% \item A language model, that takes 'prompt' from user and generates a response.
% \item Built by OpenAI
% \item Released in Nov 2022
% \item Got 1m users in 5 days  (Insta 2.5 months, Spotify 5m, Facebook 10m, Netflix 3.5yrs)
% \end{itemize}	 

% \begin{center}
% \includegraphics[width=0.6\linewidth,keepaspectratio]{chatgpt8}
% \end{center}				
% {\tiny (Ref: ChatGPT - Explained! - CodeEmporium)}
			
			
% \end{frame}

% %%%%%%%%%%%%%%%%%%%%%%%%%%%%%%%%%%%%%%%%%%%%%%%%%%%%%%%%%%%
% \begin{frame}[fragile]\frametitle{What is ChatGPT?}


% \begin{itemize}
% \item GPT based
% \item Trained on large amount of data
% \item Uses Supervised Learning and Reinforcement Learning.
% \item Gives answers like a human
% \item Can ask Follow-up questions and even admit mistakes.
% \end{itemize}	 

% \end{frame}


%%%%%%%%%%%%%%%%%%%%%%%%%%%%%%%%%%%%%%%%%%%%%%%%%%%%%%%%%%%
\begin{frame}[fragile]\frametitle{Open AI: Who are these guys?}


\begin{itemize}
\item San Francisco-based artificial intelligence company
\item Famous for its well-known DALL-E, generates images from prompts.
% \item OpenAI Inc. is the non-profit parent company of the for-profit OpenAI LP.
\item Initially supported by Elon Musk
\item The CEO is Sam Altman, who previously was president of Y Combinator.
\item Microsoft is a partner and investor.
\item Mission: To ensure that artificial general intelligence benefits all of humanity. Prevent misuse of AI
\end{itemize}	 

\begin{center}
\includegraphics[width=0.8\linewidth,keepaspectratio]{chatgpt25}
\end{center}				
{\tiny (Ref: https://www.slideegg.com/open-ai-chat-gpt)}


\end{frame}

% %%%%%%%%%%%%%%%%%%%%%%%%%%%%%%%%%%%%%%%%%%%%%%%%%%%%%%%%%%%
% \begin{frame}[fragile]\frametitle{How to Use ChatGPT?}


% \begin{center}
% \includegraphics[width=\linewidth,keepaspectratio]{chatgpt26}
% \end{center}				
% {\tiny (Ref: https://www.slideegg.com/open-ai-chat-gpt)}


% \end{frame}


% %%%%%%%%%%%%%%%%%%%%%%%%%%%%%%%%%%%%%%%%%%%%%%%%%%%%%%%%%%%
% \begin{frame}[fragile]\frametitle{GPTs Series}


% \begin{itemize}
% \item GPT + Generative Pre-trained Transformers
% \item GPT-1, GPT-2 and GPT-3 are similar in terms of architecture 
% \item Differ on the data and the number of transformer blocks with the number of incoming tokens.   
% \item GPT-1 : 12 blocks, encoded using a Byte pair encoding, 512 tokens: 117 million parameters 
% \item GPT-2 : 48 blocks, 1024 tokens: 1.5 billion parameters 
% \item GPT-3 : 96 blocks, 2048 tokens: 175 billion parameters 
% \end{itemize}	 

% \end{frame}

%%%%%%%%%%%%%%%%%%%%%%%%%%%%%%%%%%%%%%%%%%%%%%%%%%%%%%%%%%%
\begin{frame}[fragile]\frametitle{GPTs Training}

GPT: Generative Pre-trained Transformers

\begin{itemize}
% \item GPT-1 is trained in a self-supervised manner (learn to predict the next word in text data) and fine-tuned in a supervised learning manner. 
% \item GPT-2 is trained in a fully self supervised way, focusing on zero-shot transfer
% \item  GPT-3 is pre-trained in a self supervised manner exploring a bit more the few-shots fine-tuning.
\item GPT-1 is pre-trained on the BooksCorpus dataset, containing ~7000 books amounting to ~5GB of data
\item GPT-2 is pre-trained using the WebText dataset which is a more diverse set of internet data containing ~8M documents for about ~40 GB of data
\item GPT-3 uses an expanded version of the WebText dataset, two internet-based books corpora that are not disclosed and the English-language Wikipedia which constituted ~600 GB of data
\end{itemize}	 

\end{frame}

%%%%%%%%%%%%%%%%%%%%%%%%%%%%%%%%%%%%%%%%%%%%%%%%%%%%%%%%%%%
\begin{frame}[fragile]\frametitle{GPTs Training compared to human reading}

\begin{itemize}
% \item GPT-1 is trained in a self-supervised manner (learn to predict the next word in text data) and fine-tuned in a supervised learning manner. 
% \item GPT-2 is trained in a fully self supervised way, focusing on zero-shot transfer
% \item  GPT-3 is pre-trained in a self supervised manner exploring a bit more the few-shots fine-tuning.
\item GPT-3 was trained on 499B tokens; OpenAI has not published GPT-4's training size.
\item In comparison, if you spent 12 hours a day reading for an entire lifetime (80 years) at average speed (250 words / minute), we would absorb 5.26B words (tokens).
\item That's a ratio of 100:1 between the training data used for GPT-3 and the amount of data that can ever be read by a human.
\end{itemize}	 

\tiny{(Ref: LinkedIn post by Dr Jennifer Prendki)}

\end{frame}

%%%%%%%%%%%%%%%%%%%%%%%%%%%%%%%%%%%%%%%%%%%%%%%%%%%%%%%%%%%
\begin{frame}[fragile]\frametitle{GPT3 vs ChatGPT}


\begin{center}
\includegraphics[width=0.8\linewidth,keepaspectratio]{chatgpt27}
\end{center}		

\tiny{(Ref:ChatGPT Explained: Complete A-Z Guide - Kripesh Adwani)}
\end{frame}

% %%%%%%%%%%%%%%%%%%%%%%%%%%%%%%%%%%%%%%%%%%%%%%%%%%%%%%%%%%%
% \begin{frame}[fragile]\frametitle{GPT Training: Overview}


% Task: predicting next word

% \begin{center}
% \includegraphics[width=0.8\linewidth,keepaspectratio]{chatgpt13}
% \end{center}		


% {\tiny (Ref: How GPT3 Works - Easily Explained with Animations - Jay Alammar)}

% \end{frame}

% %%%%%%%%%%%%%%%%%%%%%%%%%%%%%%%%%%%%%%%%%%%%%%%%%%%%%%%%%%%
% \begin{frame}[fragile]\frametitle{GPT Inferencing: Overview}


% GPT being a decoder only model, takes words vectors (does this conversion internally), then through set of Decoder blocks, to spit out the next word one by one

% \begin{center}
% \includegraphics[width=0.8\linewidth,keepaspectratio]{chatgpt14}
% \end{center}		


% {\tiny (Ref: How GPT3 Works - Easily Explained with Animations - Jay Alammar)}

% \end{frame}

% %%%%%%%%%%%%%%%%%%%%%%%%%%%%%%%%%%%%%%%%%%%%%%%%%%%%%%%%%%%
% \begin{frame}[fragile]\frametitle{Misalignment Issue}


% \begin{itemize}
% \item GPT-3 was trained to predict the next word. Not much of context there.
% \item Introducing prompting : provide with certain samples and context

% \end{itemize}	 

			% \begin{center}
			% \includegraphics[width=0.6\linewidth,keepaspectratio]{chatgpt2}
			
			% {\tiny "Example of misalignment between text to be predicted and final use"}
			% \end{center}		
			
			
			% {\tiny (Ref: ChatGPT: training process, advantages, and limitations - By Sergio Soage, Machine Learning Engineer at Aivo)}
			
% \end{frame}

% %%%%%%%%%%%%%%%%%%%%%%%%%%%%%%%%%%%%%%%%%%%%%%%%%%%%%%%%%%%
% \begin{frame}[fragile]\frametitle{Curated Training Data}


% \begin{itemize}
% \item To fix this misalignment, humans must be involved in teaching GPT, and in this way, GPT will be able to generate better questions as it evolves.
% \item Reinforcement Learning based on Human Input (rewards if aligns) was leveraged.
% \end{itemize}	 

% \begin{center}
% \includegraphics[width=0.6\linewidth,keepaspectratio]{chatgpt3}

% \end{center}		

% {\tiny (Ref: ChatGPT: training process, advantages, and limitations - By Sergio Soage, Machine Learning Engineer at Aivo)}
			
% \end{frame}


% %%%%%%%%%%%%%%%%%%%%%%%%%%%%%%%%%%%%%%%%%%%%%%%%%%%%%%%%%%%
% \begin{frame}[fragile]\frametitle{Prerequisites: What is Reinforcement Learning (RL)?}


% \begin{itemize}
% \item RL is a method of achieving a goal via rewards. Find best path to get maximum points (game below)
% \item Agent: Makes the moves
% \item Reward: Positive or Negative (as shown)
% \item State: Current position, location.
% \item Action: Move done by the Agent
% \item Policy: Sequence of Actions that Agent takes to achieve the goal, say, Down-Down-Right-Right-Right (reward 6)..there could be other paths and one(or more) of then can be the best.
% \end{itemize}	 


% \begin{center}
% \includegraphics[width=0.6\linewidth,keepaspectratio]{chatgpt9}
% \end{center}		


% {\tiny (Ref: ChatGPT - Explained! - CodeEmporium)}
% \end{frame}

% %%%%%%%%%%%%%%%%%%%%%%%%%%%%%%%%%%%%%%%%%%%%%%%%%%%%%%%%%%%
% \begin{frame}[fragile]\frametitle{RL in context of ChatGPT}


% \begin{itemize}
% \item Agent: the model, it spits out words
% \item Reward: score at predicting each next word
% \item State: prompt + sentence generated so far
% \item Action: Predicting the next word
% \item Policy: Sequence of words and then total rewards for them.
% \item Manual person can chose between such sequences (or polices) and label which is best. Then PPO can happen to adjust backwardly the reward function.
% \item So, the Reinforcement Learning is used to capture human preferences and reduce misalignment.
% \end{itemize}	 
% \end{frame}

% %%%%%%%%%%%%%%%%%%%%%%%%%%%%%%%%%%%%%%%%%%%%%%%%%%%%%%%%%%%
% \begin{frame}[fragile]\frametitle{Better GPT 3.5}

% \begin{itemize}
% \item ChatGPT is slight deviation to InstructGPT model (GPT 3.5)
% \item InstructGPT Paper: ``Training language models to follow instructions with human feedback''
% \end{itemize}	 

% \begin{center}
% \includegraphics[width=0.8\linewidth,keepaspectratio]{chatgpt1}
% \end{center}		

% \tiny{(Ref: https://openai.com/blog/chatgpt/)}
% \end{frame}

% %%%%%%%%%%%%%%%%%%%%%%%%%%%%%%%%%%%%%%%%%%%%%%%%%%%%%%%%%%%
% \begin{frame}[fragile]\frametitle{How is ChatGPT trained?}
% Step 1

% \begin{itemize}
% \item Have a prompts dataset 
% \item Let human laborer get the correct response for them
% \item Do fine-tuning of GPT 3.5 using this data via supervised learning (SFT). 
% \item Simple but very expensive.
% \end{itemize}	 

			% \begin{center}
			% \includegraphics[width=0.8\linewidth,keepaspectratio]{chatgpt4}
			
			% \end{center}		
			
			% {\tiny (Ref: ChatGPT: training process, advantages, and limitations - By Sergio Soage, Machine Learning Engineer at Aivo)}
			
% \end{frame}

% %%%%%%%%%%%%%%%%%%%%%%%%%%%%%%%%%%%%%%%%%%%%%%%%%%%%%%%%%%%
% \begin{frame}[fragile]\frametitle{How is ChatGPT trained?}
% Step 2

% \begin{itemize}
% \item For a prompt, generate multiple responses from the fine-tuned model, one after another.
% \item Human feedback provides a ranking of responses
% \item This data is used to build Rewards Model (RM)
% \end{itemize}	 

			% \begin{center}
			% \includegraphics[width=0.8\linewidth,keepaspectratio]{chatgpt5}
			
			% \end{center}		
			
			% {\tiny (Ref: ChatGPT: training process, advantages, and limitations - By Sergio Soage, Machine Learning Engineer at Aivo)}
			
% \end{frame}


% % %%%%%%%%%%%%%%%%%%%%%%%%%%%%%%%%%%%%%%%%%%%%%%%%%%%%%%%%%%%
% % \begin{frame}[fragile]\frametitle{Step 2: Manual Ranking}


			% % \begin{center}
			% % \includegraphics[width=\linewidth,keepaspectratio]{chatgpt10}
			
			% % \end{center}		
			
			% % {\tiny (Ref: ChatGPT - Explained! - CodeEmporium)}
			
% % \end{frame}


% %%%%%%%%%%%%%%%%%%%%%%%%%%%%%%%%%%%%%%%%%%%%%%%%%%%%%%%%%%%
% \begin{frame}[fragile]\frametitle{How is ChatGPT trained?}
% Step 3

% \begin{itemize}
% \item Take unseen prompt. Pass it through Proximal Policy Optimization (PPO) model which has been initialized with SFT. It generates a response (ie inferencing via SFT).
% \item Calculate reward for that response via RM model. This reward is used to update the PPO model (loss function, back-propagation)
% \end{itemize}	 

			% \begin{center}
			% \includegraphics[width=0.8\linewidth,keepaspectratio]{chatgpt6}
			
			% \end{center}		
			
			% {\tiny (Ref: ChatGPT: training process, advantages, and limitations - By Sergio Soage, Machine Learning Engineer at Aivo)}
			
% \end{frame}

% % %%%%%%%%%%%%%%%%%%%%%%%%%%%%%%%%%%%%%%%%%%%%%%%%%%%%%%%%%%%%%%%%%%%%%%%%%%%%%%%%%%
% % \begin{frame}[fragile]\frametitle{As a Python Program?}
% % Pip install openai library, have KEY from your own account and then try various APIs available

			% % \begin{center}
			% % \includegraphics[width=0.8\linewidth,keepaspectratio]{chatgpt12}
			
			% % \end{center}		
			
			% % {\tiny (Ref: Getting Started with OpenAI API and GPT-3 -Assembly AI)}
			

% % \end{frame}


% %%%%%%%%%%%%%%%%%%%%%%%%%%%%%%%%%%%%%%%%%%%%%%%%%%%%%%%%%%%%%%%%%%%%%%%%%%%%%%%%%%
% \begin{frame}[fragile]\frametitle{Completions}

% The most common use case for any LLM is to generate a response from a prompt. In this case, the prompt is “Say this is a test”. OpenAI calls this task “completions”.

% \begin{itemize}
% \item max\_tokens: the maximum number of tokens to generate in the completion
% \item temperature: a value between 0–2. The higher the number, the more spontaneous and random the response will be.
% \item n: the number of completions to generate.
% \item stream: boolean; stream back partial progress.	
% \end{itemize}	 

% {\tiny (Ref: Techy Stuff 1: Notes on Transformers, LLMs, and OpenAI - Bill)}

% \begin{lstlisting}
% response = openai.createCompletion({
  % model="text-davinci-003",
  % prompt="Say this is a test",
  % max_tokens= 7,
  % temperature= 0,
% })
% \end{lstlisting}	


			
% \end{frame}

% %%%%%%%%%%%%%%%%%%%%%%%%%%%%%%%%%%%%%%%%%%%%%%%%%%%%%%%%%%%%%%%%%%%%%%%%%%%%%%%%%%
% \begin{frame}[fragile]\frametitle{Chat}

% A user might want to carry out a conversation with the AI. The format will be a stream of exchange between the AI (''assistant'' role) and the user (''user'' role).

% \begin{lstlisting}
% completion = openai.createChatCompletion({
  % model = "gpt-3.5-turbo",
  % messages = [{role: "user", content: "Hello world"}],
% })
% print(completion.data.choices[0].message)
% \end{lstlisting}		

			
% \end{frame}

% %%%%%%%%%%%%%%%%%%%%%%%%%%%%%%%%%%%%%%%%%%%%%%%%%%%%%%%%%%%%%%%%%%%%%%%%%%%%%%%%%%
% \begin{frame}[fragile]\frametitle{Edit}

% You might also want the AI to edit some content for you. In that case, you might want to use the edits API.

% \begin{lstlisting}
% response = openai.createEdit({
  % model= "text-davinci-edit-001",
  % input = "What day of the wek is it?",
  % instruction= "Fix the spelling mistakes",
% })
% \end{lstlisting}		

			
% \end{frame}

% %%%%%%%%%%%%%%%%%%%%%%%%%%%%%%%%%%%%%%%%%%%%%%%%%%%%%%%%%%%%%%%%%%%%%%%%%%%%%%%%%%
% \begin{frame}[fragile]\frametitle{Other Use Cases}
	

% \begin{itemize}
% \item Image: the AI will generate a new image based on your prompt and the image provided.
% \item Embeddings: turn input into a vector representation. It’s very useful when we need to compare the similarity between two texts.
% \item Audio: turn audio into text.
% \end{itemize}	 

% {\tiny (Ref: Techy Stuff 1: Notes on Transformers, LLMs, and OpenAI - Bill)}
			
% \end{frame}


\section[Apps]{Applications}
%%%%%%%%%%%%%%%%%%%%%%%%%%%%%%%%%%%%%%%%%%%%%%%%%%%%%%%%%%%%%%%%%%%%%%%%%%%%%%%%%%
\begin{frame}[fragile]\frametitle{}
\begin{center}
{\Large Applications}
\end{center}
\end{frame}

%%%%%%%%%%%%%%%%%%%%%%%%%%%%%%%%%%%%%%%%%%%%%%%%%%%%%%%%%%%%%%%%%%%%%%%%%%%%%%%%%%
\begin{frame}[fragile]\frametitle{Explain Complex Subjects}
\begin{center}
\includegraphics[width=\linewidth,keepaspectratio]{chatgpt15}
\end{center}
	
{\tiny (Ref: Top 10 Chat GPT Use Cases - Simplilearn)}
\end{frame}

%%%%%%%%%%%%%%%%%%%%%%%%%%%%%%%%%%%%%%%%%%%%%%%%%%%%%%%%%%%%%%%%%%%%%%%%%%%%%%%%%%
\begin{frame}[fragile]\frametitle{Write Code}
\begin{center}
\includegraphics[width=0.8\linewidth,keepaspectratio]{chatgpt16}
\end{center}

it was still writing but I stopped it, just to show this pic

{\tiny (Ref: Top 10 Chat GPT Use Cases - Simplilearn)}
\end{frame}

%%%%%%%%%%%%%%%%%%%%%%%%%%%%%%%%%%%%%%%%%%%%%%%%%%%%%%%%%%%%%%%%%%%%%%%%%%%%%%%%%%
\begin{frame}[fragile]\frametitle{Debug Code}
\begin{center}
\includegraphics[width=0.8\linewidth,keepaspectratio]{chatgpt17}
\end{center}

{\tiny (Ref: Top 10 Chat GPT Use Cases - Simplilearn)}
\end{frame}

%%%%%%%%%%%%%%%%%%%%%%%%%%%%%%%%%%%%%%%%%%%%%%%%%%%%%%%%%%%%%%%%%%%%%%%%%%%%%%%%%%
\begin{frame}[fragile]\frametitle{Get Custom Marketing Strategy}
\begin{center}
\includegraphics[width=\linewidth,keepaspectratio]{chatgpt18}
\end{center}

{\tiny (Ref: Top 10 Chat GPT Use Cases - Simplilearn)}
\end{frame}

%%%%%%%%%%%%%%%%%%%%%%%%%%%%%%%%%%%%%%%%%%%%%%%%%%%%%%%%%%%%%%%%%%%%%%%%%%%%%%%%%%
\begin{frame}[fragile]\frametitle{Write Articles}
\begin{center}
\includegraphics[width=\linewidth,keepaspectratio]{chatgpt19}
\end{center}

{\tiny (Ref: Top 10 Chat GPT Use Cases - Simplilearn)}
\end{frame}

%%%%%%%%%%%%%%%%%%%%%%%%%%%%%%%%%%%%%%%%%%%%%%%%%%%%%%%%%%%%%%%%%%%%%%%%%%%%%%%%%%
\begin{frame}[fragile]\frametitle{Summarize Book}
\begin{center}
\includegraphics[width=\linewidth,keepaspectratio]{chatgpt20}
\end{center}

{\tiny (Ref: Top 10 Chat GPT Use Cases - Simplilearn)}
\end{frame}

%%%%%%%%%%%%%%%%%%%%%%%%%%%%%%%%%%%%%%%%%%%%%%%%%%%%%%%%%%%%%%%%%%%%%%%%%%%%%%%%%%
\begin{frame}[fragile]\frametitle{Answer Interview Questions}
\begin{center}
\includegraphics[width=\linewidth,keepaspectratio]{chatgpt21}
\end{center}

{\tiny (Ref: Top 10 Chat GPT Use Cases - Simplilearn)}
\end{frame}



%%%%%%%%%%%%%%%%%%%%%%%%%%%%%%%%%%%%%%%%%%%%%%%%%%%%%%%%%%%%%%%%%%%%%%%%%%%%%%%%%%
\begin{frame}[fragile]\frametitle{Develop Apps}
\begin{center}
\includegraphics[width=\linewidth,keepaspectratio]{chatgpt22}
\end{center}

{\tiny (Ref: Top 10 Chat GPT Use Cases - Simplilearn)}
\end{frame}

%%%%%%%%%%%%%%%%%%%%%%%%%%%%%%%%%%%%%%%%%%%%%%%%%%%%%%%%%%%%%%%%%%%%%%%%%%%%%%%%%%
\begin{frame}[fragile]\frametitle{Create Health Plan}
\begin{center}
\includegraphics[width=\linewidth,keepaspectratio]{chatgpt23}
\end{center}

{\tiny (Ref: Top 10 Chat GPT Use Cases - Simplilearn)}
\end{frame}

%%%%%%%%%%%%%%%%%%%%%%%%%%%%%%%%%%%%%%%%%%%%%%%%%%%%%%%%%%%%%%%%%%%%%%%%%%%%%%%%%%
\begin{frame}[fragile]\frametitle{Answer General Knowledge Questions}
\begin{center}
\includegraphics[width=\linewidth,keepaspectratio]{chatgpt24}
\end{center}

{\tiny (Ref: Top 10 Chat GPT Use Cases - Simplilearn)}
\end{frame}


% %%%%%%%%%%%%%%%%%%%%%%%%%%%%%%%%%%%%%%%%%%%%%%%%%%%%%%%%%%%%%%%%%%%%%%%%%%%%%%%%%%
% \begin{frame}[fragile]\frametitle{HVAC QnA}
% Prompt: ``Diagnose a frosted evaporator coil''

% Result: 
% ``There are several reasons why an evaporator coil may become frosted:
% \begin{itemize}
% \item Low refrigerant levels: If the refrigerant levels in the system are too low, it can cause the
% evaporator coil to freeze. This is because the refrigerant is responsible for absorbing heat
% from the air, and without enough refrigerant, the coil will not be able to absorb enough heat
% to prevent freezing.
% \item Dirty air filter: A dirty air filter can restrict airflow to the evaporator coil, causing the coil
% to freeze.
% \item Blocked or restricted airflow: If the airflow to the evaporator coil is blocked or
% restricted, the coil may freeze. This can be caused by a number of things such as dirty ducts or
% registers, closed or blocked vents, or a malfunctioning blower fan.
% \item Obstruction in the drain line: A blockage in the drain line can cause water to back up and freeze on the coil.
% \end{itemize}	 
% It's important to fix the issue as soon as possible, as a frozen evaporator coil can lead to decreased
% efficiency and can cause damage to the unit.
% ''

% \tiny{(Ref:ChatGPT Could Be the 'Write' Choice for HVACR Contractors - Joanna R. Turpin)}
% \end{frame}

% %%%%%%%%%%%%%%%%%%%%%%%%%%%%%%%%%%%%%%%%%%%%%%%%%%%%%%%%%%%%%%%%%%%%%%%%%%%%%%%%%%
% \begin{frame}[fragile]\frametitle{HVAC QnA}
% Prompt: ``How to respond to a customer who thinks an air conditioner is
% too expensive''

% Result: ``
% \begin{itemize}
% \item Understand their concerns: Find out why they think the air conditioner is too expensive.
% Are they comparing it to a lower-quality unit? Are they on a tight budget? Understanding
% their concerns will help you address them more effectively.
% \item Provide a detailed explanation of the costs: Explain the costs associated with the air
% conditioner, including the cost of the unit, installation, and any additional materials or labor
% that may be required.
% \item Compare costs with the benefits: Explain the benefits of the air conditioner, such as
% energy efficiency, durability, and long-term cost savings. Explain how the cost of the unit may
% be offset by energy savings in the long run.
% \item Offer financing options: If the customer is unable to afford the air conditioner upfront,
% offer financing options that can help them spread the cost over time.

% \end{itemize}	 

% \ldots

% \tiny{(Ref:ChatGPT Could Be the 'Write' Choice for HVACR Contractors - Joanna R. Turpin)}
% \end{frame}

% %%%%%%%%%%%%%%%%%%%%%%%%%%%%%%%%%%%%%%%%%%%%%%%%%%%%%%%%%%%%%%%%%%%%%%%%%%%%%%%%%%
% \begin{frame}[fragile]\frametitle{HVAC QnA}

% \ldots 

% \begin{itemize}
% \item Provide a written estimate: Provide the customer with a detailed, written estimate that
% outlines the costs associated with the air conditioner and the benefits it will provide. This will
% help them make an informed decision.
% \item Provide a guarantee or warranty on the unit: Provide the customer with a guarantee
% or warranty on the unit to give them peace of mind in their investment.
% \item Provide an alternative solution: If the customer cannot afford the recommended
% solution, offer an alternative solution that fits their budget while still meeting their needs.
% \end{itemize}	 
% It's important to be patient, understanding and provide all the information customers need to
% make a decision that best fits their needs and budget.
% ''

% \tiny{(Ref:ChatGPT Could Be the 'Write' Choice for HVACR Contractors - Joanna R. Turpin)}
% \end{frame}

% %%%%%%%%%%%%%%%%%%%%%%%%%%%%%%%%%%%%%%%%%%%%%%%%%%%%%%%%%%%%%%%%%%%%%%%%%%%%%%%%%%
% \begin{frame}[fragile]\frametitle{If This Then That}


% \begin{itemize}
% \item With well designed prompts, one can format output in json (read, key-value pairs)
% \item Parsing via python program, one can trigger Actions. (GPT3 APIs for now)
% \item Whole pipeline can be put in place programmatically, doing any trigger based work (IFTTT - If This Then That)
% \end{itemize}	 


% \tiny{(Ref: ChatGPT for Arduino - MoonPreneur)}
% \end{frame}

% %%%%%%%%%%%%%%%%%%%%%%%%%%%%%%%%%%%%%%%%%%%%%%%%%%%%%%%%%%%%%%%%%%%%%%%%%%%%%%%%%%
% \begin{frame}[fragile]\frametitle{Possible applications}



% \begin{itemize}
% \item Multi-Language Support: ChatGPT�s ability to translate languages can be combined with Arduino to create multi-language support for devices. This feature makes it possible for users from different countries to interact with the device in their language. For example, you can create a device that provides information about local tourist attractions and supports multiple languages.
% \item Home Automation: One unique application of this combination is in the field of home automation. With ChatGPT�s ability to understand and respond to natural language and the versatility of the Arduino platform, it�s possible to create a fully voice-controlled smart home system. 
% \end{itemize}	 



% \tiny{(Ref: ChatGPT for Arduino - MoonPreneur)}
% \end{frame}


% [ai-for-x: commented 2026-11-07 | prompting is taught by ai_core_nontech_prompts (CRAFT); uncomment both lines to bring the demo back]
% \section[Demo]{Demo}
% %%%%%%%%%%%%%%%%%%%%%%%%%%%%%%%%%%%%%%%%%%%%%%%%%%%%%%%%%%%%%%%%%%%%%%%%%%%%%%%%%%
\begin{frame}[fragile]\frametitle{}
\begin{center}
{\Large Making of a Sandwich}
\end{center}
\end{frame}

%%%%%%%%%%%%%%%%%%%%%%%%%%%%%%%%%%%%%%%%%%%%%%%%%%%%%%%%%%%
% IMPROVED: Added brief explanatory note to each slide so even a layman 
% immediately understands WHY this prompt version is different/better.
% The original slides had just the prompt with no explanation.
%%%%%%%%%%%%%%%%%%%%%%%%%%%%%%%%%%%%%%%%%%%%%%%%%%%%%%%%%%%

\begin{frame}[fragile]\frametitle{Basic Prompt}

\textbf{Prompt:} \emph{Explain how to make a peanut butter and jelly sandwich.}

\medskip
\textbf{What happens:} The AI gives a generic answer, it will assume standard ingredients and a general audience. This is fine for casual use, but the output can be inconsistent if you run it multiple times.

\medskip
\textbf{Lesson:} Even a simple prompt works, but gives you \emph{no control} over the style, audience, or constraints of the answer.

\end{frame}

%%%%%%%%%%%%%%%%%%%%%%%%%%%%%%%%%%%%%%%%%%%%%%%%%%%%%%%%%%%
\begin{frame}[fragile]\frametitle{Adding Roles (Persona Prompting)}

\textbf{Prompt:} \emph{As a chef, explain to your assistant how to make a peanut butter and jelly sandwich.}

\medskip
\textbf{What changes:} The AI now adopts the voice and vocabulary of a professional chef. The tone becomes instructional and precise rather than casual.

\medskip
\textbf{Lesson:} Giving the AI a \textbf{role} (``Act as a chef / lawyer / teacher'') immediately shapes the tone, vocabulary, and depth of the response, without changing the core question.

\end{frame}

%%%%%%%%%%%%%%%%%%%%%%%%%%%%%%%%%%%%%%%%%%%%%%%%%%%%%%%%%%%
\begin{frame}[fragile]\frametitle{Adding Constraints}

\textbf{Prompt:} \emph{Make a nut-free version of the sandwich due to a customer's nut allergy.}

\medskip
\textbf{What changes:} The AI now replaces peanut butter with an alternative (e.g., sunflower seed butter or jam-only) and explicitly avoids any nut-based ingredient.

\medskip
\textbf{Lesson:} \textbf{Constraints} (``avoid X'', ``must include Y'', ``limit to N words'') are powerful control levers. The AI will respect them if you state them clearly. Vague constraints lead to vague compliance.

\end{frame}


%%%%%%%%%%%%%%%%%%%%%%%%%%%%%%%%%%%%%%%%%%%%%%%%%%%%%%%%%%%
\begin{frame}[fragile]\frametitle{Adding Examples (Few-Shot)}

\textbf{Prompt:} \emph{Create two unique variations of the classic sandwich.
Example 1, Banana Nut Crunch: [description].
Example 2, Triple Berry Blast: [description]. Now create a third variation.}

\medskip
\textbf{What changes:} By showing two examples first, you teach the AI the \textbf{format, style, and level of detail} you expect, without writing a long list of instructions.

\medskip
\textbf{Lesson:} This is called \textbf{Few-Shot Prompting}, giving the AI a few examples to pattern-match from. It is one of the most reliable ways to get consistently formatted outputs.

\end{frame}

%%%%%%%%%%%%%%%%%%%%%%%%%%%%%%%%%%%%%%%%%%%%%%%%%%%%%%%%%%%
\begin{frame}[fragile]\frametitle{Adding Contextual Information}

\textbf{Prompt:} \emph{As the head chef at ``The Sandwich Haven'', guide your new assistant to create specials for the menu.}

\medskip
\textbf{What changes:} The AI now understands the \emph{setting} (a real restaurant, not a home kitchen), the \emph{relationship} (head chef to assistant), and the \emph{purpose} (menu specials, not everyday cooking).

\medskip
\textbf{Lesson:} \textbf{Context} narrows the creative space for the AI. The richer the context, the more targeted and useful the output becomes.

\end{frame}

%%%%%%%%%%%%%%%%%%%%%%%%%%%%%%%%%%%%%%%%%%%%%%%%%%%%%%%%%%%
\begin{frame}[fragile]\frametitle{Incorporating Feedback (Iterative Prompting)}

\textbf{Prompt:} \emph{Improve the sandwich based on customer feedback: customers want less sweetness and something more creative.}

\medskip
\textbf{What changes:} The AI refines its previous output based on specific, actionable feedback rather than starting from scratch.

\medskip
\textbf{Lesson:} Prompt Engineering is an \textbf{iterative process}, you rarely get the perfect answer on the first try. Treating your conversation with the AI like a feedback loop gives progressively better results.

\end{frame}

%%%%%%%%%%%%%%%%%%%%%%%%%%%%%%%%%%%%%%%%%%%%%%%%%%%%%%%%%%%
\begin{frame}[fragile]\frametitle{Time Constraints and Prioritization}

\textbf{Prompt:} \emph{Prepare an alternative fruit-based version for testing within a tight deadline, focus on the most important steps only.}

\medskip
\textbf{What changes:} The AI prioritises and condenses, it drops elaboration and gives you a quick, actionable plan.

\medskip
\textbf{Lesson:} You can encode \textbf{business constraints} directly into your prompt: urgency, scope, priority. The AI will adapt accordingly.

\end{frame}

%%%%%%%%%%%%%%%%%%%%%%%%%%%%%%%%%%%%%%%%%%%%%%%%%%%%%%%%%%%
\begin{frame}[fragile]\frametitle{Incorporating Multidisciplinary Knowledge}

\textbf{Prompt:} \emph{Use food presentation and garnishing techniques to describe how to make a visually appealing sandwich for a food photography session.}

\medskip
\textbf{What changes:} The AI draws on knowledge from \emph{multiple domains}, culinary arts, visual presentation, and photography staging.

\medskip
\textbf{Lesson:} LLMs hold knowledge from many fields simultaneously. A good prompt can \textbf{combine domains} in ways that would take a human significant research time to replicate.

\end{frame}

%%%%%%%%%%%%%%%%%%%%%%%%%%%%%%%%%%%%%%%%%%%%%%%%%%%%%%%%%%%
\begin{frame}[fragile]\frametitle{Addressing Dietary Preferences}

\textbf{Prompt:} \emph{Prepare a fully vegan version using only plant-based alternatives for all ingredients.}

\medskip
\textbf{What changes:} The AI substitutes every non-vegan ingredient (honey, dairy butter) with appropriate alternatives and justifies each substitution.

\medskip
\textbf{Lesson:} Constraints based on \textbf{values or requirements} (vegan, halal, gluten-free, accessible, budget-friendly) are perfectly understood by modern LLMs, you just need to state them.

\end{frame}

%%%%%%%%%%%%%%%%%%%%%%%%%%%%%%%%%%%%%%%%%%%%%%%%%%%%%%%%%%%
\begin{frame}[fragile]\frametitle{Self-Criticism (Reflexion / Self-Refinement)}

\textbf{Prompt:} \emph{Explain how to make a peanut butter and jelly sandwich. Please re-read your above response. Are there any mistakes or omissions? If so, identify and correct them.}

\medskip
\textbf{What changes:} The AI reviews its own output for logical errors, missing steps, or unclear instructions, and corrects them.

\medskip
\textbf{Lesson:} You can instruct the AI to \textbf{self-critique} its answers. This technique, sometimes called \emph{Reflexion}, produces significantly more accurate and complete outputs, especially for complex tasks.

\end{frame}

%%%%%%%%%%%%%%%%%%%%%%%%%%%%%%%%%%%%%%%%%%%%%%%%%%%%%%%%%%%
\begin{frame}[fragile]\frametitle{Chain-of-Thought}

\textbf{Prompt:} \emph{Explain how to make a peanut butter and jelly sandwich. Let's think step by step.}

\medskip
\textbf{What changes:} Adding \textbf{``Let's think step by step''} triggers the AI to slow down and reason through the problem sequentially instead of jumping straight to an answer.

\medskip
\textbf{Lesson:} This tiny addition is one of the most powerful prompt engineering tricks ever discovered. It works especially well for problems that require \textbf{reasoning, planning, or multi-step logic}. The AI makes fewer mistakes when it ``shows its work''.

\end{frame}

%%%%%%%%%%%%%%%%%%%%%%%%%%%%%%%%%%%%%%%%%%%%%%%%%%%%%%%%%%%
\begin{frame}[fragile]\frametitle{Self-Consistency}

\textbf{Prompt:} \emph{Here are recipes of multiple sandwiches: Sandwich 1: [recipe 1]. Sandwich 2: [recipe 2]. \ldots\ Explain how to make a peanut butter jelly sandwich, then pick the most consistent approach across all variations.}

\medskip
\textbf{What changes:} The AI considers multiple paths to the same answer, then picks the most reliable one, like asking several people the same question and going with the majority view.

\medskip
\textbf{Lesson:} \textbf{Self-Consistency Prompting} is useful when accuracy matters more than speed. Instead of trusting one answer, you generate several reasoning paths and pick the one most answers agree on.

\end{frame}

%%%%%%%%%%%%%%%%%%%%%%%%%%%%%%%%%%%%%%%%%%%%%%%%%%%%%%%%%%%
\begin{frame}[fragile]\frametitle{Reflection and Iteration, The Big Picture}

\begin{itemize}
\item Every slide in this section illustrated one prompt technique applied to the \textbf{same task} (making a sandwich).
\item \textbf{Same AI model, same task, very different results}, only the prompt changed.
\end{itemize}

\medskip
\begin{center}
\renewcommand{\arraystretch}{1.3}
\begin{tabular}{|l|l|}
\hline
\textbf{Technique Used} & \textbf{What It Controls} \\
\hline
Role / Persona & Tone and expertise level \\
Constraints & Scope and safety \\
Few-Shot Examples & Output format and style \\
Context & Relevance and specificity \\
Feedback / Iteration & Quality improvement \\
Chain-of-Thought & Accuracy on complex tasks \\
Self-Consistency & Reliability and correctness \\
\hline
\end{tabular}
\end{center}

\medskip
\textbf{Key Takeaway:} Prompt Engineering is not magic, it is structured communication.

\end{frame}


\section[Prep]{Preparation}
%%%%%%%%%%%%%%%%%%%%%%%%%%%%%%%%%%%%%%%%%%%%%%%%%%%%%%%%%%%%%%%%%%%%%%%%%%%%%%%%%%
\begin{frame}[fragile]\frametitle{}
\begin{center}
{\Large Career in AI: But \ldots, But \ldots, How to prepare?}
\end{center}
\end{frame}


%%%%%%%%%%%%%%%%%%%%%%%%%%%%%%%%%%%%%%%%%%%%%%%%%%%%%%%%%%%
\begin{frame}[fragile]\frametitle{}
	
	\begin{center}
	{\Large What got you here,  won’t get you there!!} \\
	{\bf - Marshal Goldsmith}
	\end{center}

\end{frame}

%%%%%%%%%%%%%%%%%%%%%%%%%%%%%%%%%%%%%%%%%%%%%%%%%%%%%%%%%%%
\begin{frame}[fragile]\frametitle{}
	
	\begin{center}
	{\Large So, Well, you can’t prepare!!}  \\
	
	not everything, but certainly, specifically \ldots
	\end{center}

\end{frame}


% %%%%%%%%%%%%%%%%%%%%%%%%%%%%%%%%%%%%%%%%%%%%%%%%%%%%%%%%%%%
% \begin{frame}[fragile]\frametitle{}
	
	% \begin{center}
	% {\Large Career in Data Science \ldots}  
	% \end{center}

% \end{frame}

%%%%%%%%%%%%%%%%%%%%%%%%%%%%%%%%%%%%%%%%%%%%%%%%%%%%%%%%%%%
\begin{frame}[fragile]\frametitle{Are you suitable? }
\begin{columns}
    \begin{column}[T]{0.5\linewidth}
		Advantages:
      \begin{itemize}
			\item Domain Expertise
			\item Maturity, Communication, Soft Skills
			\item Problem Solving
			\end{itemize}
			
		Dis-advantages:
      \begin{itemize}
			\item Lost touch with Mathematics
			\item Un-Learning and Re-Learning inertia
			\item Starting from scratch? Seniority?
			\end{itemize}
			
		\end{column}
    \begin{column}[T]{0.5\linewidth}
		

			\begin{center}
			\includegraphics[width=\linewidth,keepaspectratio]{ai84}
			\end{center}		

    \end{column}
  \end{columns}
\end{frame}

%%%%%%%%%%%%%%%%%%%%%%%%%%%%%%%%%%%%%%%%%%%%%%%%%%%%%%%%%%%
\begin{frame}[fragile]\frametitle{Why do you want to Switch?}

\begin{columns}
    \begin{column}[T]{0.5\linewidth}

      \begin{itemize}
			\item $\$_1\$_2\$_3 \ldots \$_n$?
			\item Will remain in fashion forever?
			\item Hate my current job? no growth?
		\end{itemize}
			
		\end{column}
    \begin{column}[T]{0.5\linewidth}
		

			\begin{center}
			\includegraphics[width=\linewidth,keepaspectratio]{ai85}
			\end{center}		

    \end{column}
  \end{columns}
  


			
			What's in it for me?
\end{frame}

%%%%%%%%%%%%%%%%%%%%%%%%%%%%%%%%%%%%%%%%%%%%%%%%%%%%%%%%%%%
\begin{frame}[fragile]\frametitle{Current + ML combo?}

\begin{columns}
    \begin{column}[T]{0.5\linewidth}

      \begin{itemize}
			\item First : DON'T QUIT!!!
			\item Don't lose advantage due to domain expertise
			\item ML just another problem solving technique, IF DATA IS AVAILABLE
			\item Can you leverage domain expertise and apply ML there, a good/smooth transition?
			\end{itemize}
			
		\end{column}
    \begin{column}[T]{0.5\linewidth}
		

			\begin{center}
			\includegraphics[width=\linewidth,keepaspectratio]{ai86}
			\end{center}		

    \end{column}
  \end{columns}
  
  


			
\end{frame}

%%%%%%%%%%%%%%%%%%%%%%%%%%%%%%%%%%%%%%%%%%%%%%%%%%%%%%%%%%%
\begin{frame}[fragile]\frametitle{}
	
	\begin{center}
	{\Large More Generally \ldots For Career \ldots}  
	\end{center}

\end{frame}


%%%%%%%%%%%%%%%%%%%%%%%%%%%%%%%%%%%%%%%%%%%%%%%%%%%%%%%%%%%
\begin{frame}[fragile]\frametitle{Ikigai}
	
	\begin{center}
	\includegraphics[width=0.6\linewidth,keepaspectratio]{career9}
	\end{center}
	
	{\tiny (Source:  How To Find Your Ikigai And Transform Your Outlook On Life And Business - Chris Myers)}

\end{frame}

%%%%%%%%%%%%%%%%%%%%%%%%%%%%%%%%%%%%%%%%%%%%%%%%%%%%%%%%%%%
\begin{frame}[fragile]\frametitle{Generic Gyan}
\begin{columns}
    \begin{column}[T]{0.4\linewidth}
			Specific Knowledge

      \begin{itemize}
			\item Unique, rare combination
			\item Un-trainable, un-scalable
			\item Acquired through apprenticeship
			\end{itemize}
			
			Leverage
      \begin{itemize}
			\item Permission-ed: Capital, Labor
			\item Un-permission-ed: Content, code
			\item Marginal cost of duplication
			\end{itemize}
			
			Wealth is a positive sum game
			
    \end{column}
    \begin{column}[T]{0.6\linewidth}
		
			\begin{center}
			\includegraphics[width=0.4\linewidth,keepaspectratio]{career10}
			\end{center}

			Own story
      \begin{itemize}
			\item Current jobs were not available 10 yrs back
			\item Pick difficult problems/domain, Job, RnD
			\item If not now, can change later, Mid career change, harder, but possible.
			\end{itemize}
    \end{column}
  \end{columns}

{\tiny (Inspired by: Naval Ravikant, ``How to Get Rich'', 2019; twitter.com/naval)}

\end{frame}


\section[ChatGPTConcl]{ChatGPT: Conclusion}
%%%%%%%%%%%%%%%%%%%%%%%%%%%%%%%%%%%%%%%%%%%%%%%%%%%%%%%%%%%%%%%%%%%%%%%%%%%%%%%%%%
\begin{frame}[fragile]\frametitle{}
\begin{center}
{\Large Conclusions}
\end{center}
\end{frame}




%%%%%%%%%%%%%%%%%%%%%%%%%%%%%%%%%%%%%%%%%%%%%%%%%%%%%%%%%%%
\begin{frame}[fragile]\frametitle{Advantages}


\begin{itemize}
\item Conversational Abilities
\item Solving Complex Problems
\item Retaining Previous Information
% \item Creative Assistant, random-mix-ideas
\item Will replace mundane language tasks, How to articles, homework, etc
% \item Supports more complex instructions, ``reasoning'' tasks
\end{itemize}	 

\end{frame}

%%%%%%%%%%%%%%%%%%%%%%%%%%%%%%%%%%%%%%%%%%%%%%%%%%%%%%%%%%%
\begin{frame}[fragile]\frametitle{Dis-advantages}


\begin{itemize}
\item Sensitive to Input Phrasing
% \item Incorrect Answers at Times
\item Cannot replace humans for innovation, for which data does not exist already
\item Keeps ``hallucinating'': Tends to write plausible but incorrect content with confidence
\item May not get language structure right all the time, e.g try getting ghazal written
\end{itemize}	 

\end{frame}

%%%%%%%%%%%%%%%%%%%%%%%%%%%%%%%%%%%%%%%%%%%%%%%%%%%%%%%%%%%
\begin{frame}[fragile]\frametitle{Compared to humans?}


\begin{itemize}
\item Although GPT models are far more knowledgeable due to vast training (GPT-3 was trained on 499B tokens, GPT-4's size is not published, whereas a human can read about 5.26B tokens in 80 years)
\item It's not about the volume of data \ldots
\item  Children can talk long before they reach adulthood, and it's likely that $<$ 1B tokens is enough for a human to master language. Need to figure out how nature has made her neural network so efficient.
\item Imagine the pretraining done over billions of years at the time a baby comes ready with
\end{itemize}	 

\end{frame}


%%%%%%%%%%%%%%%%%%%%%%%%%%%%%%%%%%%%%%%%%%%%%%%%%%%%%%%%%%%
\begin{frame}[fragile]\frametitle{My Sketchnote}

\begin{center}
\includegraphics[width=0.45\linewidth,keepaspectratio]{ChatGPT_Sketchnote_Medium}
\end{center}		

{\tiny (Ref: https://medium.com/technology-hits/overview-of-chatgpt-95f4b43645c0)}
			

\end{frame}


%%%%%%%%%%%%%%%%%%%%%%%%%%%%%%%%%%%%%%%%%%%%%%%%%%%%%%%%%%%%%%%%%%%%%%%%%%%%%%%%%%
\begin{frame}[fragile]\frametitle{}
\begin{center}
{\Large The After-shocks}
\end{center}
\end{frame}



%%%%%%%%%%%%%%%%%%%%%%%%%%%%%%%%%%%%%%%%%%%%%%%%%%%%%%%%%%%
\begin{frame}[fragile]\frametitle{ChatGPT vs Google: Who is better?}


\begin{center}
\includegraphics[width=0.8\linewidth,keepaspectratio]{chatgpt28}
\end{center}		

Google: almost latest data, gives source, not trainable, looks more accurate
ChatGPT: a bit old data, no source, trainable, hallucinates

\tiny{(Ref:ChatGPT Explained: Complete A-Z Guide - Kripesh Adwani)}
\end{frame}


%%%%%%%%%%%%%%%%%%%%%%%%%%%%%%%%%%%%%%%%%%%%%%%%%%%%%%%%%%%
\begin{frame}[fragile]\frametitle{Will ChatGPT Kill Jobs?}


\begin{center}
\includegraphics[width=0.8\linewidth,keepaspectratio]{chatgpt29}
\end{center}		

Repetitive, boring and standard, language based jobs, for sure.
Need to be more creative, experiential to stand against ChatGPT.

\tiny{(Ref:ChatGPT Explained: Complete A-Z Guide - Kripesh Adwani)}
\end{frame}


%%%%%%%%%%%%%%%%%%%%%%%%%%%%%%%%%%%%%%%%%%%%%%%%%%%%%%%%%%%%%%%%%%%%%%%%%%%%%%%%%%
\begin{frame}[fragile]\frametitle{New Job Roles?}
Prompt Engineer: Preparing input to AI effectively to get the desired answer. Will need to AI works in the background plus domain knowledge. Give context, examples etc to prime the model to give short specific answers than the usual page-long ones (davinci GPT3 in this case)

			\begin{center}
			\includegraphics[width=0.8\linewidth,keepaspectratio]{chatgpt11}
			
			\end{center}		
			
			{\tiny (Ref: Advanced ChatGPT Guide - How to build your own Chat GPT Site - Drian Twarog)}
			

\end{frame}

% %%%%%%%%%%%%%%%%%%%%%%%%%%%%%%%%%%%%%%%%%%%%%%%%%%%%%%%%%%%
% \begin{frame}[fragile]\frametitle{Whats happening around, now?}


% \begin{itemize}
% \item ``Science journals ban listing of ChatGPT as co-author on papers'' - The Guardian
% \item OpenAI’s ChatGPT Took An IQ Test! - Two Minutes Papers
% \item Discussing internally as one of the Editors on a Medium Publication, that Do we allow AI generated articles? and if yes, how?
% \item Google’s DeepMind might release ChatGPT competitor Sparrow this year
% \end{itemize}	 

% \end{frame}

%%%%%%%%%%%%%%%%%%%%%%%%%%%%%%%%%%%%%%%%%%%%%%%%%%%%%%%%%%%
\begin{frame}[fragile]\frametitle{The Hype}


\begin{center}
\includegraphics[width=\linewidth,keepaspectratio]{chatgpt32}
\end{center}		

{\tiny (Ref: ChatGPT - Intro \& Potential Impact Sudalai Rajkumar, SRK)}
			

\end{frame}



%%%%%%%%%%%%%%%%%%%%%%%%%%%%%%%%%%%%%%%%%%%%%%%%%%%%%%%%%%%
\begin{frame}[fragile]\frametitle{Finally, from horses mouth!!}


\begin{center}
\includegraphics[width=\linewidth,keepaspectratio]{chatgpt7}
\end{center}		

{\tiny (Ref: ChatGPT: training process, advantages, and limitations - By Sergio Soage, Machine Learning Engineer at Aivo)}
			

\end{frame}

% %%%%%%%%%%%%%%%%%%%%%%%%%%%%%%%%%%%%%%%%%%%%%%%%%%%%%%%%%%%
% \begin{frame}[fragile]\frametitle{Whats next? - GPT4 }


% \begin{center}
% \includegraphics[width=\linewidth,keepaspectratio]{chatgpt33}
% \end{center}		

% *not confirmed yet

% {\tiny (Ref: https://twitter.com/WolfofBaldSt/status/1596768686051921923)}
			

% \end{frame}


%%%%%%%%%%%%%%%%%%%%%%%%%%%%%%%%%%%%%%%%%%%%%%%%%%%%%%%%%%%
\begin{frame}[fragile]\frametitle{References}
		\begin{itemize}
		\item Let's build GPT: from scratch, in code, spelled out: Andrej Karpathy
		\item ChatGPT and Reinforcement Learning - CodeEmporium
		\end{itemize}
\end{frame}


\section[Refs]{References}
%%%%%%%%%%%%%%%%%%%%%%%%%%%%%%%%%%%%%%%%%%%%%%%%%%%%%%%%%%%%%%%%%%%%%%%%%%%%%%%%%%%
\begin{frame}\frametitle{References}
Many publicly available resources have been refereed for making this presentation. Some of the notable ones are:
\small
\begin{itemize}
\item ``What counts as artificially intelligent? AI and deep learning, explained'' - The Verge
\item ``Artificial Intelligence Overview'' - Harry Surden
\item ``8 ways artificial intelligence is going to change the way you live, work and play in 2018'' - Catherine Clifford
\item ``Artificial intelligence and the future of our work'' - UNDP , IBM Research
\end{itemize}
\end{frame}

